\documentclass[11pt]{article}
\usepackage[margin=1in]{geometry}
\usepackage{amsmath,amssymb}
\usepackage{booktabs}
\usepackage{longtable}
\usepackage{array}
\usepackage[font=small,labelfont=bf]{caption}
\usepackage{microtype}
\usepackage[hidelinks]{hyperref}

\title{Reconciling the DiD-BCF Theory with the Calibration Experiments}
\author{Hugo Gobato Souto}
\date{\today}

\begin{document}
\maketitle

\begin{abstract}
\noindent The published JBES simulation grid and the new \texttt{Extra\_Theory\_Experiments}
(ETE) calibration runs share the structured DiD-BCF sampler and the same canonical
data-generating processes, so they can be compared directly. This note maps the
published estimator labels onto the new implementations, reproduces the published
numbers from the archived per-replication summaries, isolates the one pathology the
published paper flags (the doubly robust correction explodes at nonlinear degrees), and
answers two questions: whether the theory-aligned estimator improves on the old one, and
whether it remains competitive with the published benchmarks. The explosion is a
posterior \emph{tail} produced by the unregularized logistic propensity pilot under
covariate separation, not a property of the structured sampler or of the correction
algebra. At degree $d=2$ and $N=200$ two of one hundred replications carry errors $53.6$
and $25.5$ in the same-sample correction, and four of one hundred carry errors up to
$68.5$ in the reference fold convolution with the logit pilot; removing those
replications returns every RMSE to about $0.29$. The same construction with a
random-forest, intercept, or clipped pilot has no tail, and the \texttt{oracle\_canonical}
design with exact bounded propensities is stable throughout. Against the old corrected
estimator, the theory-aligned version is about the same at $d=1$ and under AR(1) errors,
and strictly better at $d=2$ and at $N=800$; against the benchmarks, the raw structured
fit remains the most accurate estimator on the common estimands, and the theory-aligned
version is competitive with Callaway--Sant'Anna and DoubleML and much better calibrated
than grf-DiD, though it does not dominate them on serial-correlation point accuracy.
The raw fit's better point accuracy is not evidence against the correction: the theory
makes no root-$n$ coverage claim for the raw summary, and the raw intervals undercover
under serial correlation at every sample size, which is exactly the failure the
correction repairs.
All numbers come from the now complete calibration archives (48 of 48 correction shards,
48 of 48 information shards, and all 384 completion shards) and the published revision summaries, through a single
reproducible script.
\end{abstract}

\section{Purpose and provenance}
\label{sec:purpose}

The JBES submission evaluates DiD-BCF with a simulation grid whose plain and corrected
summaries were produced by the revision code base
(\texttt{Simulation\_Studies\_Revision}). The ETE package reruns the same canonical and
staggered data-generating processes through the structured sampler but replaces the
single corrected estimator of that grid by a family: the raw structured posterior, the
same-sample correction reimplemented locally as \texttt{current\_hybrid}, the theorem-aligned
fixed-$K$ reference fold convolution of Algorithm~2, and the two oracle variants on an
oracle-ready design. The purpose of this note is to state precisely what matches, what
changed, and why, and to compare the new estimators with both the published DiD-BCF
numbers and the published frequentist benchmarks.

Two provenance facts matter for reading everything below. First, the simulation grid in
the submission already uses the structured sampler, so the reconciliation here concerns
the post-processing step, not the sampler; the seed instability documented for the
\texttt{mpdta} application concerns the older \texttt{published} specification and is a
separate, already documented issue. Second, the ETE runs were budgeted at the production
sampler settings ($\texttt{num\_gfr}=50$, $\texttt{num\_mcmc}=500$,
$\texttt{keep\_every}=5$, $\texttt{num\_chains}=3$), so differences from the published
grid are not budget artifacts.

Table~\ref{tab:map} maps the estimator labels used in the two code bases.

\begin{table}[htbp]
\centering
\footnotesize
\caption{Estimator map. ``Reproduced'' means the ETE run uses the same construction as
the published grid up to seeds and fold draws.}
\label{tab:map}
\begin{tabular}{p{0.20\textwidth}p{0.32\textwidth}p{0.36\textwidth}}
\toprule
Published label & ETE label & Construction \\
\midrule
DiD-BCF (plain) & \texttt{raw\_structured} & Average of raw structured posterior
$\tau$ draws over treated units in the cell. \\
DiD-BCF (corrected) & \texttt{current\_hybrid} with \texttt{logit}, \texttt{intercept},
or \texttt{rf}, optional clipping & Algorithm~1 on the same-sample posterior: Bayesian-bootstrap weights,
efficient-influence-function augmentation using a cross-fitted pilot propensity, and a
posterior bias correction using pilot $m$ and $\pi$. \\
(new) & \texttt{reference\_fold\_convolution} with \texttt{logit}, \texttt{rf}, or
\texttt{intercept} & Algorithm~2:
fixed-$K$ folds of original units, off-fold pilots, a fold-specific posterior fitted on
the evaluation fold, and a size-weighted convolution of fold draws. \\
(new) & \texttt{oracle\_current\_hybrid},
\texttt{oracle\_reference\_fold\_convolution} & The two corrections run with exact
\texttt{m0\_oracle}, \texttt{pi\_oracle}, and \texttt{barpi\_oracle} on the locally added
\texttt{oracle\_canonical} design whose propensity is Bernoulli with a known bounded
probability and whose population truth is exactly $3$. \\
\bottomrule
\end{tabular}
\end{table}

\section{Runs, estimators, and metric conventions}
\label{sec:conventions}

All errors below use replication-specific truths: for replication $r$,
$\mathrm{err}_r=\widehat\theta_r-\theta_{0,r}$, $\mathrm{SD(err)}$ is the Monte-Carlo
standard deviation of those errors, $\mathrm{RMSE}$ their root mean square,
\texttt{Cov.95} the frequency with which the reported $[q_{0.025},q_{0.975}]$ interval
contains $\theta_{0,r}$, \texttt{Len.95} its average width, and
$\overline{\mathrm{sd}}/\mathrm{SD}$ the average reported posterior standard deviation
divided by $\mathrm{SD(err)}$. The published tables divided by the raw spread of
estimates instead, so the ratio column here is \emph{recomputed} on the common error
scale; bias, $\mathrm{SD(err)}$, $\mathrm{RMSE}$, coverage, and length reproduce the
published tables exactly. Rejection rates use the two-sided $5\%$ rule that generated
the published null table, namely $\min\{P(\tau>0),P(\tau\le0)\}<0.025$.

The correction family has all $48$ shards and the information family has all $48$
shards. The completion family, which covers the cells the first two families
skipped (strong confounding, selection designs, staggered adoption, degree $d=3$,
the $N=50,100,400$ sweep cells, and the ten pre-trend designs), has all $384$
single-wave shards, so every one of its bundles has $100$ replications ($200$ for
the sharp null and the pre-trend designs); the $2{,}400$ oracle requests on the
production DGPs remain explicitly ``unavailable'' by
design, as they should, because those DGPs do not expose exact nuisances. Monte-Carlo
standard errors on a bias of magnitude $0.1$ with $\mathrm{SD(err)}\approx0.1$ are about
$0.010$, rising to $0.025$ when $\mathrm{SD(err)}=0.25$; coverage rates have Monte-Carlo
standard errors of at most $0.05$ at $p=0.5$, and of about $0.022$ at $p=0.05$. A
difference of $0.02$ to $0.03$ between an old and a new number is therefore at the edge
of resolution, and the qualitative patterns below are far larger than that.

\section{Reproduction of the published results}
\label{sec:reproduction}

Tables~\ref{tab:repro_baseline} and \ref{tab:repro_serial} compare the published grid
with the new runs for the cell that every estimator reports, $\mathrm{GATT}(4,4)$, at
$N=200$ and degrees $d=1,2$. Table~\ref{tab:repro_att} repeats the comparison for the
pooled ATT that carries the main-text table of the submission, and
Table~\ref{tab:repro800} extends the comparison to $N=800$.

The reproduction is close where it matters. At $d=1$ the new raw structured fit matches
the published plain fit closely, for instance baseline bias $-0.018$ against $-0.022$ and
$\mathrm{SD(err)}$ $0.131$ against $0.109$ (different replication panels), and the new
\texttt{current\_hybrid} matches the published corrected fit, for instance serial
coverage $0.980$ against $0.990$ and length $0.696$ against $0.693$. At $N=800$ the new
raw fit tracks the published plain fit within $0.011$ on RMSE in all four cells of
Table~\ref{tab:repro800}, and the new corrected variants match the direction and roughly
the magnitude of the published correction.

The nonlinear corner is the interesting one. At baseline $d=2$, $N=200$, the published
corrected fit has $\mathrm{RMSE}=1.040$ against a plain-fit $\mathrm{RMSE}=0.122$; the
same-sample correction in the new run has $\mathrm{RMSE}=5.947$, and the reference fold
convolution with the logit pilot has $\mathrm{RMSE}=6.894$. The raw fit is unaffected
($\mathrm{RMSE}=0.135$). The next section shows that these large numbers are carried by
a handful of replications in both runs, which is the sense in which the published
``instability at nonlinear degrees'' is reproduced rather than contradicted.

\section{The nonlinear-degree explosion is a propensity-pilot tail}
\label{sec:tail}

The per-replication errors settle the question. In the published baseline $d=2$,
$N=200$ corrected run, two of $100$ replications have $|\mathrm{err}|>2$ (the largest is
$9.41$) and deleting those two reduces RMSE from $1.040$ to $0.240$, which is the raw
scale. In the new same-sample logit run, two of $100$ replications exceed $2$ (the
largest is $53.63$) and deleting them reduces RMSE from $5.947$ to $0.294$; in the
reference fold logit run, four of $100$ exceed $2$ (the largest is $68.52$) and deleting
them reduces RMSE from $6.894$ to $0.292$. The median absolute errors are $0.159$ for the
published fit and $0.177$ to $0.190$ for the two new logit fits, so the central behavior
of all three is the same and only the tail differs.

The diagnostics identify the cause. In the same-sample logit run at that cell the
largest observed control odds is at the clipping bound, $\widehat\pi/(1-\widehat\pi)
=999$, the effective sample size falls to $1.02$, and the maximum normalized control
weight reaches $0.988$; in the reference fold logit run the maximum odds are $986$ with
minimum effective sample size $1.15$. The $d=2$ covariates are Kang--Schafer transforms
including $e^{X_4/2}$ and cubic terms, so an unregularized logistic fit separates the
groups, the inverse-odds factor in the augmentation and in $\widehat\gamma$ becomes
enormous, and a single control observation can dominate the corrected draw. The failure
is a violation of the bounded-overlap and localization conditions of the BvM theorem
(Conditions on allocation and rates in Section S7 of the theory supplement require
$1-\widehat\pi$ and $\widehat\gamma$ to stay in fixed $L^\infty$ balls), so the observed
tail lies outside the theorem rather than contradicting it.

Three pieces of evidence confirm the diagnosis. First, capping the odds at $19$ by
setting the clipping constant to $0.05$ eliminates the tail: in the same cell the
maximum $|\mathrm{err}|$ falls from $53.63$ to $0.770$, RMSE falls from $5.947$ to
$0.258$, coverage is $0.980$, and the average length is $1.062$. Second, replacing the
logit pilot by the random-forest pilot, whose odds range from a median of $2.5$ to a
maximum of $5.3$ with minimum effective sample size $33.1$, removes the tail without any
clipping: maximum $|\mathrm{err}|=0.839$ and $\mathrm{RMSE}=0.234$, with coverage
$0.970$ and length $0.986$. The intercept pilot behaves similarly
($\mathrm{RMSE}=0.214$). Third, Table~\ref{tab:oracle} reports the
\texttt{oracle\_canonical} design, in which the propensity is known to lie in
$(0.001,0.999)$ and $m_0$ is exact. There the $d=2$, $N=200$ corrected variants all stay
between $\mathrm{RMSE}=0.199$ (current RF) and $0.265$ (reference logit), with no tail
at all, and at $N=800$ every corrected variant has $\mathrm{RMSE}$ between $0.109$ and
$0.130$. Holding the correction fixed and improving overlap removes the pathology;
holding overlap fixed and changing the fold construction does not.

\section{Performance of the reference fold convolution}
\label{sec:reference}

The reference fold convolution is the estimator the theory actually covers, so its
finite-sample profile deserves a direct reading. At $d=1$, $N=200$, baseline, its
random-forest version has bias $-0.030$, $\mathrm{SD(err)}=0.243$,
$\mathrm{RMSE}=0.243$, coverage $0.910$, and length $0.962$, against published corrected
bias $-0.028$, $\mathrm{SD(err)}=0.231$, $\mathrm{RMSE}=0.231$, coverage $0.950$,
length $0.979$. It is therefore about as efficient as the published correction at that
cell, with slightly wider intervals and coverage about $1.8$ Monte-Carlo standard errors
below nominal; its value is that it is the construction the BvM theorem describes, and
with a stable pilot it remains controlled at $d=2$ where the old published run did not
($\mathrm{RMSE}=0.234$ against $1.040$). At $N=800$ the reference variants track the
same-sample variants closely, for example baseline $d=1$ reference RF
$\mathrm{RMSE}=0.130$ against published corrected $0.107$ and new current logit $0.109$,
with coverage $0.920$ against $0.980$ and $0.940$ and length $0.477$ against $0.450$.

Two honest qualifications follow. The correction, in any of its implementations, pays
an efficiency cost at $N=200$: in the baseline cells the raw $\mathrm{RMSE}$ is about
$0.13$ against $0.21$ to $0.26$ for the corrected variants, with intervals roughly
$1.6$ to $1.9$ times longer, although under serial correlation at $d=2$ the corrected
variants are close to or better than the raw fit ($0.145$ to $0.164$ against $0.141$).
And the reference construction does not by itself fix the pilot: with the logit pilot it
inherits the tail described in Section~\ref{sec:tail}. The practical recommendation is
to pair the theorem-aligned fold construction with a pilot that respects overlap, which
the random-forest pilot does out of the box in these runs, or to report clipping
explicitly as an exploratory stabilization.

\section{Completion grid: the cells the first runs skipped}
\label{sec:completion}

Tables~\ref{tab:completion_b1}, \ref{tab:completion_sweep},
\ref{tab:completion_staggered}, \ref{tab:completion_null},
\ref{tab:completion_pt}, and \ref{tab:completion_pt_d3} report the completion
family: the strong-confounder, selection, staggered, degree $d=3$, small-$N$
sweep, sharp-null $d=3$, and pre-trend cells that the correction-audit family
did not cover, each with the new raw fit, the new reference-fold fit, and the
published same-sample correction alongside for contrast. The pattern from the
old cells carries over, with two honest qualifications. First, where the
same-sample correction already had no tail, the reference fold pays a modest
efficiency cost for the fold split: at baseline $d=1$, $N=400$, the published
corrected RMSE is $0.154$ against $0.189$ for the reference, and under serial
correlation $0.090$ against $0.125$; at $N=50$, where the old run did have a
tail, the reference wins ($0.532$ against $0.707$ at baseline $d=1$), while at
$N=100$ the two are close ($0.410$ against $0.342$). The reference
removes tails; it does not dominate point estimation where no tail existed.
Second, the reference-fold pre-trend diagnostic is oversized:
its slope rule rejects $0.105$ of the time at $d=1$ under \texttt{PT\_hold}
($0.135$ at $d=3$) against $0.015$ ($0.040$) for the raw diagnostic, so the raw
diagnostic remains the testing instrument and the fold version an exploratory
sensitivity.

Strong confounding is now a clean win. At $d=1$ the reference ATT has bias
$0.005$, RMSE $0.184$, coverage $0.98$, side by side with the published
$0.008$, $0.176$, $0.99$; at $d=2$, $-0.009$, $0.221$, $0.95$,
against $0.037$, $0.359$, $0.96$ for the published correction; at $d=3$,
$0.033$, $0.192$, $0.97$ against $0.033$, $0.361$, $0.97$. Selection on
observables behaves the same way at $d=2$ and $d=3$ ($0.240$ and $0.259$
against $1.338$ and $1.361$), while at $d=1$, where the same-sample fit never
had a tail, the two are side by side ($0.232$ against $0.244$). Selection on
both observables and unobservables, a scenario the paper never ran, is
controlled at every degree ($0.195$, $0.244$, $0.258$ with coverage $0.96$,
$0.92$, $0.94$). The $d=3$ baseline and serial cells, where the old run blew up
to RMSE $1.096$ and $0.601$, are now $0.219$ and $0.198$ with coverage $0.94$
and $0.90$.

Staggered adoption is controlled at every degree. The reference event-study
biases at $d=1$ are $-0.046$, $-0.023$, $-0.001$, $-0.031$ across $k=0,\dots,3$
with coverage $0.96$ to $0.98$; the nine cohort-time cells hold coverage $0.96$
to $1.00$ with biases under $0.061$ in absolute value, including the late
cohorts where the raw fit degrades. At $d=2$ and $d=3$ the same picture holds
(Table~\ref{tab:completion_staggered}): the $3.455$ and $3.461$ cell-averaged
RMSEs of the old $d=2,3$ runs become $0.366$ and $0.365$.

The sharp-null $d=3$ rerun (Table~\ref{tab:completion_null}, $200$
replications) is the one place where the reference fit overshoots: it rejects
$0.130$ of ATT tests, $0.100$ of event-study tests, and $0.100$ of cohort-time
tests at nominal $5\%$, against $0.050$, $0.030$, $0.030$ for the raw fit. With
$200$ replications the Monte-Carlo standard error at $p=0.13$ is $0.024$, so
$0.13$ is about three standard errors above nominal: a genuine mild
oversizing, not noise, and it is reported as such in the paper. Coverage tells
the same story ($0.87$, $0.90$, $0.90$). It remains far from the grf-DiD
distortion, and the $d=1,2$ null cells stay at $0.04$ to $0.09$.

The pre-trend diagnostic comparison (Tables~\ref{tab:completion_pt} and
\ref{tab:completion_pt_d3}) confirms the raw diagnostic and bounds the fold
version. The raw slope rule reproduces the published operating characteristics
on fresh replication panels (hold size $0.015$ at $d=1$, power $0.130$,
$0.385$, $0.855$, $1.000$ at slopes $0.05$ through $0.40$), and the $X_1$
contrast carries the heterogeneous violations ($1.000$ at $\kappa=0.4$ for both
versions) while staying quiet elsewhere. The fold version detects more at every
magnitude but pays for it in size, as noted above, and its slope intervals
undercover ($0.89$ at hold, $0.73$ under \texttt{PT\_conditional}). The
discrimination property belongs to the raw diagnostic: it rejects $0.060$
under \texttt{PT\_conditional} where the marginal event study rejects $0.755$,
while the fold version rejects $0.275$.

\section{Is the theory-aligned estimator better than the old corrected one?}
\label{sec:oldnew}

Table~\ref{tab:headtohead} puts the two corrections side by side in the cells where
they matter. The answer is that it is about the same in the well-behaved cells and
strictly better in the cells that used to fail, which is the same as saying that the
theory-aligned construction removes the failure without giving up what worked.

At $d=1$ and $N=200$, baseline, the published corrected RMSE is $0.231$ with coverage
$0.950$ and the new reference RF RMSE is $0.243$ with coverage $0.910$; under serial
correlation at $d=1$ the numbers are $0.136$ with coverage $0.990$ against $0.160$ with
coverage $0.960$. Those differences are small and within about one to two Monte-Carlo
standard errors, with the reference variant slightly more conservative in length and
slightly lower in coverage. The same-sample logit variant is closer to the published
correction in these cells ($0.214$ and $0.145$), as expected, since it is the same
algebra.

At $d=2$ the comparison stops being close. At $N=200$, baseline, the published corrected
RMSE is $1.040$; the reference RF is $0.234$ with coverage $0.970$, the same-sample RF is
$0.224$ with coverage $0.970$, and the clipped logit is $0.258$ with coverage $0.980$.
At $N=800$, baseline, the published corrected RMSE is $0.318$ against $0.114$ for the
reference RF and $0.102$ for the same-sample logit. Under serial correlation at $d=2$,
$N=200$, the published corrected RMSE is $0.269$ against $0.164$ for the reference RF.

The verdict is therefore: same performance in the benign regimes, clearly better in the
nonlinear and large-$N$ regimes, and no longer at the mercy of a handful of replications.
The one caveat is that the reference construction inherits whatever the pilot does, so
the improvement comes from the joint choice of the off-fold Algorithm~2 construction and
an overlap-respecting pilot, not from Algorithm~2 alone.

\section{Is the theory-aligned estimator still better than the benchmarks?}
\label{sec:benchmarks}

Tables~\ref{tab:benchgatt}, \ref{tab:benchgatt800}, and \ref{tab:benchatt} compare the
new estimators with the published frequentist benchmarks on the same designs and
replication-specific truths, and Table~\ref{tab:benchhead} condenses the
$\mathrm{GATT}(4,4)$ comparison. The short answer is that the raw structured fit is
still the accuracy leader, while the theory-aligned corrected fit is now competitive
with the best benchmark rather than dominated by it.

On point accuracy at $N=200$, baseline, $d=1$, the raw fit has $\mathrm{RMSE}=0.131$
against $0.221$ for Callaway--Sant'Anna, $0.304$ for DoubleML, and $0.222$ for grf-DiD.
The same-sample correction has $\mathrm{RMSE}=0.214$, matching Callaway--Sant'Anna and
beating DoubleML, and the reference RF has $0.243$, between Callaway--Sant'Anna and
DoubleML. At $d=2$ the raw fit is again first ($0.135$), and the corrected versions are
$0.224$ to $0.258$ against $0.225$ for Callaway--Sant'Anna, $0.280$ for DoubleML, and
$0.223$ for grf-DiD, so they are level with the two best benchmarks. Coverage favors
the new corrected versions over grf-DiD by a wide margin ($0.97$ against $0.81$ at both
degrees) and is comparable to Callaway--Sant'Anna and DoubleML.

The one place where the corrected fit does not match the best benchmark is serial
correlation. At $d=1$, $N=200$, the raw fit has the best $\mathrm{RMSE}$ ($0.126$) but
coverage $0.920$, Callaway--Sant'Anna has $0.129$ with coverage $0.990$, and the
reference RF has $0.160$ with coverage $0.960$; at $d=2$ the pattern repeats with
Callaway--Sant'Anna at $0.135$ against $0.164$ for the reference RF. So the corrected
DiD-BCF pays a visible efficiency cost on serially correlated data relative to the
purpose-built efficient benchmark, while remaining close to DoubleML and clearly better
calibrated than grf-DiD. At $N=800$ the raw fit dominates every benchmark by a wide
margin ($\mathrm{RMSE}=0.069$ against $0.107$ to $0.114$), and the corrected versions
are at $0.109$ to $0.130$, in the same neighborhood as the benchmarks.

The pooled ATT tells the same story (Table~\ref{tab:benchatt}). At baseline $d=1$ the
raw fit has $\mathrm{RMSE}=0.112$ against $0.179$ for TWFE and $0.169$ for Synthetic
DiD; the same-sample correction has $0.176$, essentially TWFE, and the reference RF has
$0.203$, a little worse than both but with coverage $0.930$. Under serial correlation
at $d=1$ the raw fit's coverage drops to $0.810$ while the same-sample correction holds
$0.980$ at $\mathrm{RMSE}=0.142$, better than TWFE ($0.198$) and close to Synthetic DiD
($0.166$), which is exactly the trade the theory predicts.

In one sentence: the raw structured DiD-BCF remains the most accurate estimator on every
common cell, the theory-aligned corrected version is now competitive with
Callaway--Sant'Anna and DoubleML and far better calibrated than grf-DiD, and it no
longer collapses at nonlinear degrees, but it does not uniformly dominate
Callaway--Sant'Anna under serial correlation, where the efficient frequentist benchmark
is still the sharper instrument for a point estimate.

\section{Does the raw fit's accuracy indicate that the correction is wrong?}
\label{sec:rawtruth}

No, and the experiments show why. ``Raw'' means the plain structured posterior summary,
the average of the $\tau$ draws over treated units in the cell, before any
post-processing. It is the estimator the paper recommends for heterogeneity analyses and
point summaries, and the theory deliberately makes no root-$n$ coverage claim for it:
Section S7 states that uncorrected intervals require an additional equicontinuity regime
and that raw plug-in summaries carry no root-$n$ guarantee, and the minimax lower bound
in the main text makes the complementary point that no estimator can have a uniform
root-$n$ rate for the CATT surface, so any average of it inherits either a
smoothness-limited rate or an extra localization requirement.

Three features of the numbers confirm that the correction, not the raw fit, is doing its
job. First, the raw intervals are the ones that fail when they should: under AR(1)
errors their coverage is $0.88$ to $0.92$ at $N=200$ and remains $0.89$ to $0.92$ at
$N=800$, a missing variance component that does not vanish with sample size, while the
corrected variants cover at $0.96$ to $1.00$ at both sizes. A point estimator can have a
low RMSE and still be uncalibrated; here it is. Second, at $N=800$ the corrected
interval lengths match the asymptotically efficient benchmark almost exactly (baseline
$d=1$: $0.451$ against $0.445$ for Callaway--Sant'Anna; serial $d=1$: $0.329$ against
$0.326$), which is the quantitative content of the BvM theorem; the raw intervals are
much shorter ($0.286$ and $0.234$) precisely because they are not calibrated to the
sampling variability of the efficient average. Third, the oracle design, where the
nuisance conditions of the theorem hold exactly, produces corrected RMSEs of $0.199$ to
$0.265$ at $d=2$, $N=200$ with coverage $0.95$ to $0.99$, and no tail; if the
correction algebra were wrong, exact nuisances would not rescue it.

The raw fit's better RMSE at $N=200$ is a bias-variance outcome. The effect-forest prior
shrinks toward smooth, small effects, and these DGPs have smooth low-dimensional effect
surfaces, so at small $N$ the shrinkage wins while the correction pays the variance of
estimating the nuisance pilots and the augmentation. That is a finite-sample property of
a well-matched prior and a smooth truth, not a failure of an asymptotic statement, and
the efficiency claim is a limiting local statement about the corrected law. The two live
side by side: report the raw fit for point accuracy and heterogeneity, and the corrected
fit for calibrated inference, which is exactly the paper's recommendation.

\section{Information ablation}
\label{sec:information}

Tables~\ref{tab:infogatt} and~\ref{tab:infoatt} report the information ablation.
The reduced-cell fit uses only cohort $g$ and never-treated units at periods $g-1$ and
$t$; the pooled fit keeps the full panel but forces a single effect surface with
\texttt{effect\_by\_cohort=False}.

For the canonical designs the reduction is viable but costly. At baseline $d=1$,
$N=800$, the reduced cell raises $\mathrm{GATT}(4,4)$ RMSE from $0.065$ to $0.111$ and
roughly doubles the interval length from $0.288$ to $0.493$, with coverage rising from
$0.95$ to $0.98$; at $N=200$ the RMSE rises from $0.127$ to $0.207$ with length from
$0.559$ to $0.936$. Under serial correlation the reduction is nearly free in RMSE at
$N=800$ ($0.060$ against $0.058$) and slightly favorable at $N=200$ $d=2$ ($0.126$
against $0.134$), always with coverage at or above nominal. The consistent pattern is
that discarding the other periods and cohorts costs some efficiency and a good deal of
length, not validity.

Pooling is a different matter. With a single treated cohort, baseline and serial, the
pooled fit is numerically almost identical to the full-panel fit, for example baseline
$d=1$, $N=800$, RMSE $0.067$ against $0.065$ and coverage $0.97$ against $0.95$, because
there is nothing to pool across. With staggered cohorts it is catastrophic: the pooled
fit reports a single surface for cohorts whose effects differ, and the
$\mathrm{GATT}(4,4)$ bias is $1.417$ at $N=200$ and $1.467$ at $N=800$ with coverage
$0.00$ in every cell, while the reduced-cell fit stays approximately unbiased
($0.012$ and $0.022$) with coverage $1.00$. The pooled ATT is similarly biased
($-0.204$ at $N=800$ with coverage $0.13$). The ablation therefore says that the
two-group two-period reduction is a defensible local device and that pooling cohorts is
not, which is the expected reading of the no-already-treated-as-controls principle.

\section{Sharp-null size}
\label{sec:null}

Table~\ref{tab:null} reports the two-sided $5\%$ rejection rate under the sharp null.
With the complete data the picture is cleaner than in the partial run. The raw
structured fit is conservative everywhere, between $0.010$ and $0.080$. The published
corrected numbers in the matched cells are $0.040$ to $0.090$, the new same-sample
variants are between $0.000$ and $0.070$, and the reference variants are between $0.040$
and $0.100$. Every cell is within about $1.7$ Monte-Carlo standard errors of nominal,
including the largest cells, the reference-logit \texttt{ATT} entries at $0.100$, whose
Monte-Carlo standard error is $0.030$. The correction therefore neither restores exact
size nor destroys it: it trades some of the raw fit's conservativeness for better
interval calibration, at nominal size.

The $d=3$ sharp-null rerun ($200$ replications, Table~\ref{tab:completion_null})
is the one cell where that trade goes too far. The raw fit stays conservative
(rejection $0.050$, $0.030$, $0.030$ on ATT, event-study, and cohort-time cells),
while the reference fold rejects $0.130$, $0.100$, $0.100$, with coverage $0.87$,
$0.90$, $0.90$. At $n=200$ the Monte-Carlo standard error at $p=0.13$ is $0.024$,
so the ATT rejection rate is about three standard errors above nominal: a genuine
mild oversizing at $d=3$, reported as such, not hidden inside the $d=1,2$ range.
It is still far from the grf-DiD distortion, and it does not overturn the
calibration story elsewhere, but a reader who needs exact size at $d=3$ should
not rely on the corrected intervals alone.

\section{What the reconciliation establishes}
\label{sec:conclusion}

With all correction shards in hand the account is complete, including the
completion grid that the first runs skipped. The published plain and
corrected central behavior is reproduced with the current structured sampler; the
published instability of the corrected estimator at $d\ge2$ is reproduced and localized
to a small-replication tail caused by inverse odds under logistic separation; the tail
exists equally in the theorem-aligned reference construction when that construction is
handed the same pilot, and disappears with the random-forest, intercept, or clipped
pilots, and by construction on the oracle design. The completion rerun extends the
same verdict to strong confounding, selection designs, staggered adoption, degree
$d=3$, and the small-$N$ sweep cells, with the two documented qualifications: a
modest efficiency cost where the old run had no tail, and a mildly oversized
$d=3$ sharp-null test and an oversized fold-pooled pre-trend diagnostic. The theory-aligned construction fixes
the independence and weighting structure of the correction, while the overlap conditions
remain an empirical obligation on the pilot. The paper's advice to report both fits side
by side remains correct, with the amendment that the corrected fit should name its pilot
and that a logit pilot on transformed covariates with $10^{-3}$ clipping should not be
the default.

The experiments here do not establish any theorem. They do not cover the contracted
posterior rate, they use finite
posterior and bootstrap draw budgets, and the oracle design is deliberately idealized
(bounded propensity and iid row errors). They are calibration evidence about which
implementation of the theory behaves as advertised.

\section{Limitations}
\label{sec:limitations}

Most calibration cells have $100$ replications; the completion sharp-null and
pre-trend reruns use $200$. The published and new runs use
different replication panels, so the exact magnitudes of tail-driven quantities are
replication-specific; this is why Section~\ref{sec:tail} emphasizes the diagnostic and
the trimmed RMSE rather than the raw RMSE. The benchmark comparison in
Section~\ref{sec:benchmarks} uses the published benchmark runs on the same DGP designs,
so it is a comparison across independent panels rather than a paired comparison. The
staggered comparison between the published \texttt{D\_staggered} run and the new
information family is consistent in direction and roughly two Monte-Carlo standard
errors apart in bias, which we read as agreement within seed allocation. Finally,
clipping is an exploratory stabilization and is not theoretically neutral, so the
clipped results should be read as a sensitivity check alongside the random-forest pilot,
not as a theorem-covered default.

\clearpage
{\footnotesize
\begin{longtable}{lrrrrrrr}
\caption{$\mathrm{GATT}(g{=}4,t{=}4)$, $N=200$, baseline design: published JBES run (100 replications) against the Extra\_Theory calibration run (available replications shown). Metrics use replication-specific truths; here $\overline{\mathrm{sd}}/\mathrm{SD}$ is average reported posterior SD over the Monte-Carlo error SD $\mathrm{SD}(\widehat\theta-\theta_0)$, which is the common scale used throughout this note (the published tables divided by the raw spread of estimates instead).}\label{tab:repro_baseline}\\
\toprule
Estimator & Bias & SD(err) & RMSE & Cov.95 & Len.95 & $\overline{\mathrm{sd}}/\mathrm{SD}$ & $n$ \\
\midrule

\endfirsthead
\toprule
Estimator & Bias & SD(err) & RMSE & Cov.95 & Len.95 & $\overline{\mathrm{sd}}/\mathrm{SD}$ & $n$ \\
\midrule

\endhead
\bottomrule
\endlastfoot
\multicolumn{8}{l}{\textit{degree $d=1$}} \\
\quad JBES plain (structured) & -0.022 & 0.109 & 0.110 & 1.000 & 0.558 & 1.305 & 100 \\
\quad JBES corrected (same-sample) & -0.028 & 0.231 & 0.231 & 0.950 & 0.979 & 1.082 & 100 \\
\quad New raw structured & -0.018 & 0.131 & 0.131 & 0.980 & 0.561 & 1.097 & 100 \\
\quad New current hybrid, logit & -0.036 & 0.212 & 0.214 & 0.960 & 0.965 & 1.160 & 100 \\
\quad New current hybrid, logit clip .05 & -0.036 & 0.212 & 0.214 & 0.960 & 0.965 & 1.160 & 100 \\
\quad New current hybrid, RF & -0.029 & 0.208 & 0.209 & 0.980 & 0.937 & 1.146 & 100 \\
\quad New reference fold, logit & -0.031 & 0.254 & 0.254 & 0.910 & 0.994 & 0.998 & 100 \\
\quad New reference fold, RF & -0.030 & 0.243 & 0.243 & 0.910 & 0.962 & 1.010 & 100 \\
\quad New reference fold, intercept & -0.029 & 0.221 & 0.222 & 0.940 & 0.897 & 1.036 & 100 \\
\addlinespace
\multicolumn{8}{l}{\textit{degree $d=2$}} \\
\quad JBES plain (structured) & -0.024 & 0.120 & 0.122 & 0.980 & 0.572 & 1.214 & 100 \\
\quad JBES corrected (same-sample) & 0.109 & 1.039 & 1.040 & 0.960 & 1.676 & 0.409 & 100 \\
\quad New raw structured & 0.018 & 0.134 & 0.135 & 0.960 & 0.575 & 1.090 & 100 \\
\quad New current hybrid, logit & 0.796 & 5.923 & 5.947 & 0.980 & 4.024 & 0.179 & 100 \\
\quad New current hybrid, logit clip .05 & 0.021 & 0.259 & 0.258 & 0.980 & 1.062 & 1.046 & 100 \\
\quad New current hybrid, RF & 0.020 & 0.224 & 0.224 & 0.970 & 0.965 & 1.097 & 100 \\
\quad New reference fold, logit & 0.651 & 6.898 & 6.894 & 0.920 & 3.546 & 0.141 & 100 \\
\quad New reference fold, RF & 0.022 & 0.234 & 0.234 & 0.970 & 0.986 & 1.073 & 100 \\
\quad New reference fold, intercept & 0.022 & 0.214 & 0.214 & 0.960 & 0.910 & 1.087 & 100 \\
\addlinespace
\end{longtable}
}

{\footnotesize
\begin{longtable}{lrrrrrrr}
\caption{$\mathrm{GATT}(g{=}4,t{=}4)$, $N=200$, AR(1) serial design: published JBES run (100 replications) against the Extra\_Theory calibration run (available replications shown). Metrics use replication-specific truths; here $\overline{\mathrm{sd}}/\mathrm{SD}$ is average reported posterior SD over the Monte-Carlo error SD $\mathrm{SD}(\widehat\theta-\theta_0)$, which is the common scale used throughout this note (the published tables divided by the raw spread of estimates instead).}\label{tab:repro_serial}\\
\toprule
Estimator & Bias & SD(err) & RMSE & Cov.95 & Len.95 & $\overline{\mathrm{sd}}/\mathrm{SD}$ & $n$ \\
\midrule

\endfirsthead
\toprule
Estimator & Bias & SD(err) & RMSE & Cov.95 & Len.95 & $\overline{\mathrm{sd}}/\mathrm{SD}$ & $n$ \\
\midrule

\endhead
\bottomrule
\endlastfoot
\multicolumn{8}{l}{\textit{degree $d=1$}} \\
\quad JBES plain (structured) & -0.012 & 0.111 & 0.111 & 0.970 & 0.462 & 1.058 & 100 \\
\quad JBES corrected (same-sample) & -0.024 & 0.134 & 0.136 & 0.990 & 0.693 & 1.317 & 100 \\
\quad New raw structured & 0.003 & 0.127 & 0.126 & 0.920 & 0.465 & 0.935 & 100 \\
\quad New current hybrid, logit & -0.008 & 0.146 & 0.145 & 0.980 & 0.696 & 1.219 & 100 \\
\quad New current hybrid, logit clip .05 & -0.008 & 0.146 & 0.145 & 0.980 & 0.696 & 1.219 & 100 \\
\quad New current hybrid, RF & -0.009 & 0.144 & 0.144 & 0.980 & 0.672 & 1.192 & 100 \\
\quad New reference fold, logit & -0.003 & 0.185 & 0.184 & 0.940 & 0.717 & 0.987 & 100 \\
\quad New reference fold, RF & -0.012 & 0.161 & 0.160 & 0.960 & 0.687 & 1.090 & 100 \\
\quad New reference fold, intercept & -0.010 & 0.154 & 0.154 & 0.950 & 0.646 & 1.066 & 100 \\
\addlinespace
\multicolumn{8}{l}{\textit{degree $d=2$}} \\
\quad JBES plain (structured) & -0.013 & 0.127 & 0.128 & 0.960 & 0.478 & 0.957 & 100 \\
\quad JBES corrected (same-sample) & 0.019 & 0.269 & 0.269 & 0.990 & 1.160 & 1.101 & 100 \\
\quad New raw structured & -0.016 & 0.141 & 0.141 & 0.880 & 0.480 & 0.871 & 100 \\
\quad New current hybrid, logit & -0.025 & 0.264 & 0.264 & 0.990 & 0.969 & 0.935 & 100 \\
\quad New current hybrid, logit clip .05 & -0.001 & 0.159 & 0.158 & 0.990 & 0.761 & 1.223 & 100 \\
\quad New current hybrid, RF & 0.002 & 0.145 & 0.145 & 1.000 & 0.688 & 1.207 & 100 \\
\quad New reference fold, logit & -0.014 & 0.327 & 0.325 & 0.950 & 0.942 & 0.738 & 100 \\
\quad New reference fold, RF & 0.003 & 0.165 & 0.164 & 0.980 & 0.701 & 1.083 & 100 \\
\quad New reference fold, intercept & 0.003 & 0.147 & 0.147 & 0.990 & 0.655 & 1.137 & 100 \\
\addlinespace
\end{longtable}
}

{\footnotesize
\begin{longtable}{lrrrrrrr}
\caption{Pooled ATT, $N=200$: published JBES run against the new calibration runs. The ratio column is recomputed on the common error-SD scale.}\label{tab:repro_att}\\
\toprule
Estimator & Bias & SD(err) & RMSE & Cov.95 & Len.95 & $\overline{\mathrm{sd}}/\mathrm{SD}$ & $n$ \\
\midrule

\endfirsthead
\toprule
Estimator & Bias & SD(err) & RMSE & Cov.95 & Len.95 & $\overline{\mathrm{sd}}/\mathrm{SD}$ & $n$ \\
\midrule

\endhead
\bottomrule
\endlastfoot
\multicolumn{8}{l}{\textit{baseline, degree $d=1$}} \\
\quad JBES plain (structured) & 0.000 & 0.098 & 0.098 & 0.960 & 0.425 & 1.105 & 100 \\
\quad JBES corrected (same-sample) & -0.000 & 0.180 & 0.179 & 0.980 & 0.813 & 1.152 & 100 \\
\quad New raw structured & -0.007 & 0.113 & 0.112 & 0.950 & 0.425 & 0.967 & 100 \\
\quad New current hybrid, logit & -0.015 & 0.177 & 0.176 & 0.950 & 0.794 & 1.148 & 100 \\
\quad New current hybrid, logit clip .05 & -0.015 & 0.177 & 0.176 & 0.950 & 0.794 & 1.148 & 100 \\
\quad New current hybrid, RF & -0.014 & 0.172 & 0.171 & 0.960 & 0.773 & 1.148 & 100 \\
\quad New reference fold, logit & -0.028 & 0.213 & 0.213 & 0.930 & 0.821 & 0.982 & 100 \\
\quad New reference fold, RF & -0.015 & 0.203 & 0.203 & 0.930 & 0.793 & 0.994 & 100 \\
\quad New reference fold, intercept & -0.006 & 0.195 & 0.194 & 0.960 & 0.746 & 0.977 & 100 \\
\addlinespace
\multicolumn{8}{l}{\textit{baseline, degree $d=2$}} \\
\quad JBES plain (structured) & -0.003 & 0.116 & 0.115 & 0.940 & 0.437 & 0.968 & 100 \\
\quad JBES corrected (same-sample) & 0.116 & 1.069 & 1.070 & 0.970 & 1.502 & 0.357 & 100 \\
\quad New raw structured & 0.014 & 0.120 & 0.120 & 0.950 & 0.442 & 0.943 & 100 \\
\quad New current hybrid, logit & 0.468 & 4.008 & 4.015 & 0.990 & 2.838 & 0.186 & 100 \\
\quad New current hybrid, logit clip .05 & 0.004 & 0.215 & 0.214 & 0.990 & 0.900 & 1.068 & 100 \\
\quad New current hybrid, RF & 0.008 & 0.172 & 0.171 & 0.980 & 0.794 & 1.179 & 100 \\
\quad New reference fold, logit & 0.579 & 5.759 & 5.759 & 0.870 & 2.812 & 0.132 & 100 \\
\quad New reference fold, RF & 0.022 & 0.180 & 0.180 & 0.970 & 0.820 & 1.161 & 100 \\
\quad New reference fold, intercept & 0.024 & 0.161 & 0.162 & 0.970 & 0.761 & 1.209 & 100 \\
\addlinespace
\multicolumn{8}{l}{\textit{serial, degree $d=1$}} \\
\quad JBES plain (structured) & 0.005 & 0.127 & 0.126 & 0.830 & 0.360 & 0.732 & 100 \\
\quad JBES corrected (same-sample) & -0.002 & 0.150 & 0.150 & 0.980 & 0.704 & 1.193 & 100 \\
\quad New raw structured & 0.004 & 0.130 & 0.130 & 0.810 & 0.361 & 0.714 & 100 \\
\quad New current hybrid, logit & -0.010 & 0.143 & 0.142 & 0.980 & 0.692 & 1.237 & 100 \\
\quad New current hybrid, logit clip .05 & -0.010 & 0.143 & 0.142 & 0.980 & 0.692 & 1.237 & 100 \\
\quad New current hybrid, RF & -0.012 & 0.143 & 0.143 & 0.970 & 0.674 & 1.201 & 100 \\
\quad New reference fold, logit & -0.003 & 0.190 & 0.189 & 0.940 & 0.710 & 0.952 & 100 \\
\quad New reference fold, RF & -0.018 & 0.172 & 0.172 & 0.930 & 0.680 & 1.010 & 100 \\
\quad New reference fold, intercept & -0.013 & 0.166 & 0.166 & 0.940 & 0.641 & 0.985 & 100 \\
\addlinespace
\multicolumn{8}{l}{\textit{serial, degree $d=2$}} \\
\quad JBES plain (structured) & 0.000 & 0.144 & 0.143 & 0.820 & 0.378 & 0.672 & 100 \\
\quad JBES corrected (same-sample) & 0.071 & 0.477 & 0.480 & 0.960 & 1.234 & 0.640 & 100 \\
\quad New raw structured & -0.018 & 0.149 & 0.149 & 0.770 & 0.381 & 0.657 & 100 \\
\quad New current hybrid, logit & -0.003 & 0.332 & 0.330 & 0.980 & 1.091 & 0.836 & 100 \\
\quad New current hybrid, logit clip .05 & -0.011 & 0.179 & 0.179 & 0.980 & 0.746 & 1.064 & 100 \\
\quad New current hybrid, RF & -0.009 & 0.161 & 0.160 & 0.980 & 0.698 & 1.108 & 100 \\
\quad New reference fold, logit & -0.052 & 0.558 & 0.558 & 0.880 & 0.954 & 0.434 & 100 \\
\quad New reference fold, RF & -0.017 & 0.197 & 0.196 & 0.940 & 0.711 & 0.922 & 100 \\
\quad New reference fold, intercept & -0.010 & 0.171 & 0.171 & 0.950 & 0.660 & 0.988 & 100 \\
\addlinespace
\end{longtable}
}

{\footnotesize
\begin{longtable}{lrrrrrrr}
\caption{$\mathrm{GATT}(g{=}4,t{=}4)$, $N=800$: the published $B2$ sweep against the new calibration runs.}\label{tab:repro800}\\
\toprule
Estimator & Bias & SD(err) & RMSE & Cov.95 & Len.95 & $\overline{\mathrm{sd}}/\mathrm{SD}$ & $n$ \\
\midrule

\endfirsthead
\toprule
Estimator & Bias & SD(err) & RMSE & Cov.95 & Len.95 & $\overline{\mathrm{sd}}/\mathrm{SD}$ & $n$ \\
\midrule

\endhead
\bottomrule
\endlastfoot
\multicolumn{8}{l}{\textit{baseline, degree $d=1$}} \\
\quad JBES plain (structured) & 0.002 & 0.069 & 0.068 & 0.970 & 0.288 & 1.073 & 100 \\
\quad JBES corrected (same-sample) & 0.001 & 0.107 & 0.107 & 0.980 & 0.450 & 1.072 & 100 \\
\quad New raw structured & 0.008 & 0.069 & 0.069 & 0.960 & 0.286 & 1.059 & 100 \\
\quad New current hybrid, logit & 0.004 & 0.109 & 0.109 & 0.940 & 0.451 & 1.058 & 100 \\
\quad New current hybrid, logit clip .05 & 0.004 & 0.109 & 0.109 & 0.940 & 0.451 & 1.058 & 100 \\
\quad New current hybrid, RF & 0.001 & 0.110 & 0.109 & 0.950 & 0.471 & 1.098 & 100 \\
\quad New reference fold, logit & -0.001 & 0.122 & 0.122 & 0.920 & 0.456 & 0.954 & 100 \\
\quad New reference fold, RF & -0.004 & 0.130 & 0.130 & 0.920 & 0.477 & 0.936 & 100 \\
\quad New reference fold, intercept & 0.000 & 0.114 & 0.113 & 0.930 & 0.446 & 1.002 & 100 \\
\addlinespace
\multicolumn{8}{l}{\textit{baseline, degree $d=2$}} \\
\quad JBES plain (structured) & -0.001 & 0.073 & 0.072 & 0.960 & 0.302 & 1.061 & 100 \\
\quad JBES corrected (same-sample) & -0.046 & 0.316 & 0.318 & 0.990 & 0.624 & 0.511 & 100 \\
\quad New raw structured & 0.009 & 0.060 & 0.061 & 0.990 & 0.301 & 1.273 & 100 \\
\quad New current hybrid, logit & 0.005 & 0.103 & 0.102 & 0.960 & 0.453 & 1.129 & 100 \\
\quad New current hybrid, logit clip .05 & 0.005 & 0.103 & 0.102 & 0.960 & 0.453 & 1.129 & 100 \\
\quad New current hybrid, RF & 0.003 & 0.100 & 0.099 & 0.960 & 0.476 & 1.223 & 100 \\
\quad New reference fold, logit & 0.008 & 0.109 & 0.108 & 0.950 & 0.457 & 1.078 & 100 \\
\quad New reference fold, RF & 0.007 & 0.114 & 0.114 & 0.950 & 0.479 & 1.073 & 100 \\
\quad New reference fold, intercept & 0.006 & 0.103 & 0.102 & 0.960 & 0.446 & 1.111 & 100 \\
\addlinespace
\multicolumn{8}{l}{\textit{serial, degree $d=1$}} \\
\quad JBES plain (structured) & 0.003 & 0.067 & 0.067 & 0.900 & 0.235 & 0.896 & 100 \\
\quad JBES corrected (same-sample) & 0.006 & 0.071 & 0.071 & 0.990 & 0.328 & 1.177 & 100 \\
\quad New raw structured & 0.001 & 0.063 & 0.063 & 0.920 & 0.234 & 0.959 & 100 \\
\quad New current hybrid, logit & 0.003 & 0.065 & 0.064 & 1.000 & 0.329 & 1.300 & 100 \\
\quad New current hybrid, logit clip .05 & 0.003 & 0.065 & 0.064 & 1.000 & 0.329 & 1.300 & 100 \\
\quad New current hybrid, RF & 0.002 & 0.067 & 0.067 & 1.000 & 0.341 & 1.289 & 100 \\
\quad New reference fold, logit & 0.001 & 0.068 & 0.068 & 1.000 & 0.330 & 1.235 & 100 \\
\quad New reference fold, RF & 0.002 & 0.072 & 0.072 & 0.980 & 0.343 & 1.211 & 100 \\
\quad New reference fold, intercept & 0.002 & 0.066 & 0.065 & 1.000 & 0.324 & 1.262 & 100 \\
\addlinespace
\multicolumn{8}{l}{\textit{serial, degree $d=2$}} \\
\quad JBES plain (structured) & 0.000 & 0.071 & 0.071 & 0.910 & 0.249 & 0.898 & 100 \\
\quad JBES corrected (same-sample) & -0.014 & 0.176 & 0.176 & 1.000 & 0.431 & 0.634 & 100 \\
\quad New raw structured & 0.001 & 0.074 & 0.074 & 0.910 & 0.252 & 0.872 & 100 \\
\quad New current hybrid, logit & -0.005 & 0.065 & 0.065 & 0.980 & 0.337 & 1.332 & 100 \\
\quad New current hybrid, logit clip .05 & -0.005 & 0.065 & 0.065 & 0.980 & 0.335 & 1.322 & 100 \\
\quad New current hybrid, RF & -0.006 & 0.064 & 0.064 & 0.990 & 0.346 & 1.391 & 100 \\
\quad New reference fold, logit & -0.007 & 0.077 & 0.077 & 0.980 & 0.346 & 1.151 & 100 \\
\quad New reference fold, RF & -0.004 & 0.074 & 0.074 & 0.990 & 0.347 & 1.201 & 100 \\
\quad New reference fold, intercept & -0.006 & 0.062 & 0.062 & 0.990 & 0.327 & 1.337 & 100 \\
\addlinespace
\end{longtable}
}

{\footnotesize
\begin{longtable}{lrrrrrrr}
\caption{\texttt{oracle\_canonical} design (known bounded $\pi(X)\in(0.001,0.999)$, exact $m_0$; population truth $=3$): the new estimators only.}\label{tab:oracle}\\
\toprule
Estimator & Bias & SD(err) & RMSE & Cov.95 & Len.95 & $\overline{\mathrm{sd}}/\mathrm{SD}$ & $n$ \\
\midrule

\endfirsthead
\toprule
Estimator & Bias & SD(err) & RMSE & Cov.95 & Len.95 & $\overline{\mathrm{sd}}/\mathrm{SD}$ & $n$ \\
\midrule

\endhead
\bottomrule
\endlastfoot
\multicolumn{8}{l}{\textit{$N=200$, degree $d=1$}} \\
\quad New raw structured & 0.008 & 0.112 & 0.112 & 1.000 & 0.526 & 1.199 & 100 \\
\quad New current hybrid, logit & 0.004 & 0.260 & 0.259 & 0.910 & 0.988 & 0.961 & 100 \\
\quad New current hybrid, RF & 0.004 & 0.236 & 0.235 & 0.910 & 0.858 & 0.924 & 100 \\
\quad New reference fold, logit & 0.033 & 0.323 & 0.323 & 0.890 & 1.028 & 0.807 & 100 \\
\quad New reference fold, RF & 0.016 & 0.279 & 0.278 & 0.900 & 0.890 & 0.811 & 100 \\
\quad New oracle current hybrid & 0.008 & 0.248 & 0.247 & 0.930 & 0.934 & 0.954 & 100 \\
\quad New oracle reference fold & 0.009 & 0.254 & 0.253 & 0.930 & 0.947 & 0.946 & 100 \\
\addlinespace
\multicolumn{8}{l}{\textit{$N=200$, degree $d=2$}} \\
\quad New raw structured & 0.009 & 0.121 & 0.121 & 0.970 & 0.519 & 1.088 & 100 \\
\quad New current hybrid, logit & -0.025 & 0.222 & 0.222 & 0.960 & 1.020 & 1.163 & 100 \\
\quad New current hybrid, RF & -0.011 & 0.200 & 0.199 & 0.950 & 0.862 & 1.099 & 100 \\
\quad New reference fold, logit & -0.044 & 0.262 & 0.265 & 0.950 & 1.054 & 1.017 & 100 \\
\quad New reference fold, RF & -0.025 & 0.239 & 0.239 & 0.910 & 0.896 & 0.952 & 100 \\
\quad New oracle current hybrid & -0.007 & 0.221 & 0.220 & 0.970 & 0.985 & 1.130 & 100 \\
\quad New oracle reference fold & -0.008 & 0.224 & 0.223 & 0.970 & 1.001 & 1.128 & 100 \\
\addlinespace
\multicolumn{8}{l}{\textit{$N=800$, degree $d=1$}} \\
\quad New raw structured & 0.005 & 0.066 & 0.066 & 0.970 & 0.278 & 1.067 & 100 \\
\quad New current hybrid, logit & -0.007 & 0.144 & 0.143 & 0.920 & 0.478 & 0.849 & 100 \\
\quad New current hybrid, RF & -0.002 & 0.146 & 0.145 & 0.920 & 0.484 & 0.847 & 100 \\
\quad New reference fold, logit & -0.001 & 0.157 & 0.157 & 0.880 & 0.485 & 0.786 & 100 \\
\quad New reference fold, RF & -0.006 & 0.161 & 0.160 & 0.890 & 0.495 & 0.786 & 100 \\
\quad New oracle current hybrid & -0.000 & 0.137 & 0.136 & 0.940 & 0.470 & 0.878 & 100 \\
\quad New oracle reference fold & -0.000 & 0.137 & 0.136 & 0.940 & 0.472 & 0.879 & 100 \\
\addlinespace
\multicolumn{8}{l}{\textit{$N=800$, degree $d=2$}} \\
\quad New raw structured & 0.016 & 0.064 & 0.066 & 0.970 & 0.282 & 1.124 & 100 \\
\quad New current hybrid, logit & 0.014 & 0.117 & 0.117 & 0.970 & 0.479 & 1.047 & 100 \\
\quad New current hybrid, RF & 0.016 & 0.120 & 0.120 & 0.970 & 0.494 & 1.050 & 100 \\
\quad New reference fold, logit & 0.015 & 0.122 & 0.123 & 0.960 & 0.481 & 1.005 & 100 \\
\quad New reference fold, RF & 0.014 & 0.129 & 0.130 & 0.930 & 0.493 & 0.972 & 100 \\
\quad New oracle current hybrid & 0.011 & 0.109 & 0.109 & 0.970 & 0.476 & 1.118 & 100 \\
\quad New oracle reference fold & 0.011 & 0.109 & 0.109 & 0.970 & 0.477 & 1.119 & 100 \\
\addlinespace
\end{longtable}
}

{\footnotesize
\begin{longtable}{lrrrrr}
\caption{$\mathrm{GATT}(g{=}4,t{=}4)$: the new DiD-BCF calibration estimators against the published JBES benchmarks (100 replications per cell, replication-specific truths).}\label{tab:benchgatt}\\
\toprule
Estimator & Bias & SD(err) & RMSE & Cov.95 & Len.95 \\
\midrule
\endfirsthead
\toprule
Estimator & Bias & SD(err) & RMSE & Cov.95 & Len.95 \\
\midrule
\endhead
\bottomrule
\endlastfoot
\multicolumn{6}{l}{\textit{baseline, $N=200$, degree $d=1$}} \\
\quad JBES plain (structured) & -0.022 & 0.109 & 0.110 & 1.000 & 0.558 \\
\quad JBES corrected (same-sample) & -0.028 & 0.231 & 0.231 & 0.950 & 0.979 \\
\quad DiD-BCF raw (new) & -0.018 & 0.131 & 0.131 & 0.980 & 0.561 \\
\quad DiD-BCF same-sample, logit (new) & -0.036 & 0.212 & 0.214 & 0.960 & 0.965 \\
\quad DiD-BCF same-sample, logit clip .05 (new) & -0.036 & 0.212 & 0.214 & 0.960 & 0.965 \\
\quad DiD-BCF reference fold, RF (new) & -0.030 & 0.243 & 0.243 & 0.910 & 0.962 \\
\quad DiD-BCF reference fold, logit (new) & -0.031 & 0.254 & 0.254 & 0.910 & 0.994 \\
\quad Callaway--Sant'Anna & -0.031 & 0.220 & 0.221 & 0.970 & 0.902 \\
\quad DoubleML DR-DiD & -0.039 & 0.303 & 0.304 & 0.920 & 1.157 \\
\quad grf-DiD (Wang) & -0.036 & 0.220 & 0.222 & 0.810 & 0.630 \\
\addlinespace
\multicolumn{6}{l}{\textit{baseline, $N=200$, degree $d=2$}} \\
\quad JBES plain (structured) & -0.024 & 0.120 & 0.122 & 0.980 & 0.572 \\
\quad JBES corrected (same-sample) & 0.109 & 1.039 & 1.040 & 0.960 & 1.676 \\
\quad DiD-BCF raw (new) & 0.018 & 0.134 & 0.135 & 0.960 & 0.575 \\
\quad DiD-BCF same-sample, logit (new) & 0.796 & 5.923 & 5.947 & 0.980 & 4.024 \\
\quad DiD-BCF same-sample, logit clip .05 (new) & 0.021 & 0.259 & 0.258 & 0.980 & 1.062 \\
\quad DiD-BCF reference fold, RF (new) & 0.022 & 0.234 & 0.234 & 0.970 & 0.986 \\
\quad DiD-BCF reference fold, logit (new) & 0.651 & 6.898 & 6.894 & 0.920 & 3.546 \\
\quad Callaway--Sant'Anna & -0.034 & 0.224 & 0.225 & 0.960 & 0.914 \\
\quad DoubleML DR-DiD & -0.034 & 0.279 & 0.280 & 0.940 & 1.127 \\
\quad grf-DiD (Wang) & -0.037 & 0.221 & 0.223 & 0.810 & 0.630 \\
\addlinespace
\multicolumn{6}{l}{\textit{serial, $N=200$, degree $d=1$}} \\
\quad JBES plain (structured) & -0.012 & 0.111 & 0.111 & 0.970 & 0.462 \\
\quad JBES corrected (same-sample) & -0.024 & 0.134 & 0.136 & 0.990 & 0.693 \\
\quad DiD-BCF raw (new) & 0.003 & 0.127 & 0.126 & 0.920 & 0.465 \\
\quad DiD-BCF same-sample, logit (new) & -0.008 & 0.146 & 0.145 & 0.980 & 0.696 \\
\quad DiD-BCF same-sample, logit clip .05 (new) & -0.008 & 0.146 & 0.145 & 0.980 & 0.696 \\
\quad DiD-BCF reference fold, RF (new) & -0.012 & 0.161 & 0.160 & 0.960 & 0.687 \\
\quad DiD-BCF reference fold, logit (new) & -0.003 & 0.185 & 0.184 & 0.940 & 0.717 \\
\quad Callaway--Sant'Anna & -0.025 & 0.127 & 0.129 & 0.990 & 0.653 \\
\quad DoubleML DR-DiD & -0.029 & 0.163 & 0.165 & 0.990 & 0.776 \\
\quad grf-DiD (Wang) & -0.032 & 0.130 & 0.133 & 0.910 & 0.458 \\
\addlinespace
\multicolumn{6}{l}{\textit{serial, $N=200$, degree $d=2$}} \\
\quad JBES plain (structured) & -0.013 & 0.127 & 0.128 & 0.960 & 0.478 \\
\quad JBES corrected (same-sample) & 0.019 & 0.269 & 0.269 & 0.990 & 1.160 \\
\quad DiD-BCF raw (new) & -0.016 & 0.141 & 0.141 & 0.880 & 0.480 \\
\quad DiD-BCF same-sample, logit (new) & -0.025 & 0.264 & 0.264 & 0.990 & 0.969 \\
\quad DiD-BCF same-sample, logit clip .05 (new) & -0.001 & 0.159 & 0.158 & 0.990 & 0.761 \\
\quad DiD-BCF reference fold, RF (new) & 0.003 & 0.165 & 0.164 & 0.980 & 0.701 \\
\quad DiD-BCF reference fold, logit (new) & -0.014 & 0.327 & 0.325 & 0.950 & 0.942 \\
\quad Callaway--Sant'Anna & -0.028 & 0.133 & 0.135 & 0.990 & 0.662 \\
\quad DoubleML DR-DiD & -0.022 & 0.170 & 0.170 & 0.970 & 0.773 \\
\quad grf-DiD (Wang) & -0.032 & 0.130 & 0.133 & 0.900 & 0.458 \\
\addlinespace
\end{longtable}
}

{\footnotesize
\begin{longtable}{lrrrrr}
\caption{$\mathrm{GATT}(g{=}4,t{=}4)$: the new DiD-BCF calibration estimators against the published JBES benchmarks (100 replications per cell, replication-specific truths).}\label{tab:benchgatt800}\\
\toprule
Estimator & Bias & SD(err) & RMSE & Cov.95 & Len.95 \\
\midrule
\endfirsthead
\toprule
Estimator & Bias & SD(err) & RMSE & Cov.95 & Len.95 \\
\midrule
\endhead
\bottomrule
\endlastfoot
\multicolumn{6}{l}{\textit{baseline, $N=800$, degree $d=1$}} \\
\quad JBES plain (structured) & 0.002 & 0.069 & 0.068 & 0.970 & 0.288 \\
\quad JBES corrected (same-sample) & 0.001 & 0.107 & 0.107 & 0.980 & 0.450 \\
\quad DiD-BCF raw (new) & 0.008 & 0.069 & 0.069 & 0.960 & 0.286 \\
\quad DiD-BCF same-sample, logit (new) & 0.004 & 0.109 & 0.109 & 0.940 & 0.451 \\
\quad DiD-BCF same-sample, logit clip .05 (new) & 0.004 & 0.109 & 0.109 & 0.940 & 0.451 \\
\quad DiD-BCF reference fold, RF (new) & -0.004 & 0.130 & 0.130 & 0.920 & 0.477 \\
\quad DiD-BCF reference fold, logit (new) & -0.001 & 0.122 & 0.122 & 0.920 & 0.456 \\
\quad Callaway--Sant'Anna & 0.002 & 0.108 & 0.107 & 0.980 & 0.445 \\
\quad DoubleML DR-DiD & 0.005 & 0.115 & 0.114 & 0.970 & 0.503 \\
\quad grf-DiD (Wang) & 0.001 & 0.112 & 0.111 & 0.810 & 0.319 \\
\addlinespace
\multicolumn{6}{l}{\textit{baseline, $N=800$, degree $d=2$}} \\
\quad JBES plain (structured) & -0.001 & 0.073 & 0.072 & 0.960 & 0.302 \\
\quad JBES corrected (same-sample) & -0.046 & 0.316 & 0.318 & 0.990 & 0.624 \\
\quad DiD-BCF raw (new) & 0.009 & 0.060 & 0.061 & 0.990 & 0.301 \\
\quad DiD-BCF same-sample, logit (new) & 0.005 & 0.103 & 0.102 & 0.960 & 0.453 \\
\quad DiD-BCF same-sample, logit clip .05 (new) & 0.005 & 0.103 & 0.102 & 0.960 & 0.453 \\
\quad DiD-BCF reference fold, RF (new) & 0.007 & 0.114 & 0.114 & 0.950 & 0.479 \\
\quad DiD-BCF reference fold, logit (new) & 0.008 & 0.109 & 0.108 & 0.950 & 0.457 \\
\quad Callaway--Sant'Anna & 0.001 & 0.108 & 0.108 & 0.980 & 0.449 \\
\quad DoubleML DR-DiD & -0.000 & 0.113 & 0.113 & 0.990 & 0.515 \\
\quad grf-DiD (Wang) & 0.001 & 0.112 & 0.111 & 0.810 & 0.319 \\
\addlinespace
\multicolumn{6}{l}{\textit{serial, $N=800$, degree $d=1$}} \\
\quad JBES plain (structured) & 0.003 & 0.067 & 0.067 & 0.900 & 0.235 \\
\quad JBES corrected (same-sample) & 0.006 & 0.071 & 0.071 & 0.990 & 0.328 \\
\quad DiD-BCF raw (new) & 0.001 & 0.063 & 0.063 & 0.920 & 0.234 \\
\quad DiD-BCF same-sample, logit (new) & 0.003 & 0.065 & 0.064 & 1.000 & 0.329 \\
\quad DiD-BCF same-sample, logit clip .05 (new) & 0.003 & 0.065 & 0.064 & 1.000 & 0.329 \\
\quad DiD-BCF reference fold, RF (new) & 0.002 & 0.072 & 0.072 & 0.980 & 0.343 \\
\quad DiD-BCF reference fold, logit (new) & 0.001 & 0.068 & 0.068 & 1.000 & 0.330 \\
\quad Callaway--Sant'Anna & 0.005 & 0.070 & 0.070 & 0.990 & 0.326 \\
\quad DoubleML DR-DiD & 0.006 & 0.082 & 0.082 & 0.950 & 0.365 \\
\quad grf-DiD (Wang) & 0.007 & 0.074 & 0.074 & 0.890 & 0.244 \\
\addlinespace
\multicolumn{6}{l}{\textit{serial, $N=800$, degree $d=2$}} \\
\quad JBES plain (structured) & 0.000 & 0.071 & 0.071 & 0.910 & 0.249 \\
\quad JBES corrected (same-sample) & -0.014 & 0.176 & 0.176 & 1.000 & 0.431 \\
\quad DiD-BCF raw (new) & 0.001 & 0.074 & 0.074 & 0.910 & 0.252 \\
\quad DiD-BCF same-sample, logit (new) & -0.005 & 0.065 & 0.065 & 0.980 & 0.337 \\
\quad DiD-BCF same-sample, logit clip .05 (new) & -0.005 & 0.065 & 0.065 & 0.980 & 0.335 \\
\quad DiD-BCF reference fold, RF (new) & -0.004 & 0.074 & 0.074 & 0.990 & 0.347 \\
\quad DiD-BCF reference fold, logit (new) & -0.007 & 0.077 & 0.077 & 0.980 & 0.346 \\
\quad Callaway--Sant'Anna & 0.004 & 0.071 & 0.071 & 0.980 & 0.329 \\
\quad DoubleML DR-DiD & -0.004 & 0.076 & 0.076 & 1.000 & 0.366 \\
\quad grf-DiD (Wang) & 0.007 & 0.074 & 0.074 & 0.880 & 0.244 \\
\addlinespace
\end{longtable}
}

{\footnotesize
\begin{longtable}{lrrrrr}
\caption{Pooled ATT: the new DiD-BCF calibration estimators against the published JBES benchmarks (100 replications per cell, replication-specific truths).}\label{tab:benchatt}\\
\toprule
Estimator & Bias & SD(err) & RMSE & Cov.95 & Len.95 \\
\midrule
\endfirsthead
\toprule
Estimator & Bias & SD(err) & RMSE & Cov.95 & Len.95 \\
\midrule
\endhead
\bottomrule
\endlastfoot
\multicolumn{6}{l}{\textit{baseline, $N=200$, degree $d=1$}} \\
\quad JBES plain (structured) & 0.000 & 0.098 & 0.098 & 0.960 & 0.425 \\
\quad JBES corrected (same-sample) & -0.000 & 0.180 & 0.179 & 0.980 & 0.813 \\
\quad DiD-BCF raw (new) & -0.007 & 0.113 & 0.112 & 0.950 & 0.425 \\
\quad DiD-BCF same-sample, logit (new) & -0.015 & 0.177 & 0.176 & 0.950 & 0.794 \\
\quad DiD-BCF same-sample, logit clip .05 (new) & -0.015 & 0.177 & 0.176 & 0.950 & 0.794 \\
\quad DiD-BCF reference fold, RF (new) & -0.015 & 0.203 & 0.203 & 0.930 & 0.793 \\
\quad DiD-BCF reference fold, logit (new) & -0.028 & 0.213 & 0.213 & 0.930 & 0.821 \\
\quad TWFE & -0.024 & 0.178 & 0.179 & 0.980 & 0.837 \\
\quad Synthetic DiD & -0.023 & 0.168 & 0.169 & 1.000 & 0.797 \\
\addlinespace
\multicolumn{6}{l}{\textit{baseline, $N=200$, degree $d=2$}} \\
\quad JBES plain (structured) & -0.003 & 0.116 & 0.115 & 0.940 & 0.437 \\
\quad JBES corrected (same-sample) & 0.116 & 1.069 & 1.070 & 0.970 & 1.502 \\
\quad DiD-BCF raw (new) & 0.014 & 0.120 & 0.120 & 0.950 & 0.442 \\
\quad DiD-BCF same-sample, logit (new) & 0.468 & 4.008 & 4.015 & 0.990 & 2.838 \\
\quad DiD-BCF same-sample, logit clip .05 (new) & 0.004 & 0.215 & 0.214 & 0.990 & 0.900 \\
\quad DiD-BCF reference fold, RF (new) & 0.022 & 0.180 & 0.180 & 0.970 & 0.820 \\
\quad DiD-BCF reference fold, logit (new) & 0.579 & 5.759 & 5.759 & 0.870 & 2.812 \\
\quad TWFE & -0.024 & 0.178 & 0.179 & 0.980 & 0.837 \\
\quad Synthetic DiD & -0.023 & 0.168 & 0.169 & 1.000 & 0.797 \\
\addlinespace
\multicolumn{6}{l}{\textit{serial, $N=200$, degree $d=1$}} \\
\quad JBES plain (structured) & 0.005 & 0.127 & 0.126 & 0.830 & 0.360 \\
\quad JBES corrected (same-sample) & -0.002 & 0.150 & 0.150 & 0.980 & 0.704 \\
\quad DiD-BCF raw (new) & 0.004 & 0.130 & 0.130 & 0.810 & 0.361 \\
\quad DiD-BCF same-sample, logit (new) & -0.010 & 0.143 & 0.142 & 0.980 & 0.692 \\
\quad DiD-BCF same-sample, logit clip .05 (new) & -0.010 & 0.143 & 0.142 & 0.980 & 0.692 \\
\quad DiD-BCF reference fold, RF (new) & -0.018 & 0.172 & 0.172 & 0.930 & 0.680 \\
\quad DiD-BCF reference fold, logit (new) & -0.003 & 0.190 & 0.189 & 0.940 & 0.710 \\
\quad TWFE & -0.022 & 0.198 & 0.198 & 0.980 & 0.894 \\
\quad Synthetic DiD & -0.024 & 0.165 & 0.166 & 0.990 & 0.792 \\
\addlinespace
\multicolumn{6}{l}{\textit{serial, $N=200$, degree $d=2$}} \\
\quad JBES plain (structured) & 0.000 & 0.144 & 0.143 & 0.820 & 0.378 \\
\quad JBES corrected (same-sample) & 0.071 & 0.477 & 0.480 & 0.960 & 1.234 \\
\quad DiD-BCF raw (new) & -0.018 & 0.149 & 0.149 & 0.770 & 0.381 \\
\quad DiD-BCF same-sample, logit (new) & -0.003 & 0.332 & 0.330 & 0.980 & 1.091 \\
\quad DiD-BCF same-sample, logit clip .05 (new) & -0.011 & 0.179 & 0.179 & 0.980 & 0.746 \\
\quad DiD-BCF reference fold, RF (new) & -0.017 & 0.197 & 0.196 & 0.940 & 0.711 \\
\quad DiD-BCF reference fold, logit (new) & -0.052 & 0.558 & 0.558 & 0.880 & 0.954 \\
\quad TWFE & -0.022 & 0.198 & 0.198 & 0.980 & 0.894 \\
\quad Synthetic DiD & -0.024 & 0.165 & 0.166 & 0.990 & 0.792 \\
\addlinespace
\end{longtable}
}

{\footnotesize
\begin{longtable}{llrrrrrr}
\caption{$\mathrm{GATT}(g{=}4,t{=}4)$ in the information-ablation family (100 replications per cell). ``reduced'' uses only the two groups and periods of the cell; ``pooled'' fits one effect surface with \texttt{effect\_by\_cohort=False}.}\label{tab:infogatt}\\
\toprule
Design & $d$ & $N$ & Variant & Bias & RMSE & Cov.95 & Len.95 \\
\midrule
\endfirsthead
\toprule
Design & $d$ & $N$ & Variant & Bias & RMSE & Cov.95 & Len.95 \\
\midrule
\endhead
\bottomrule
\endlastfoot
baseline & 1 & 200 & full panel & 0.005 & 0.127 & 0.970 & 0.559 \\
baseline & 1 & 200 & reduced cell & 0.018 & 0.207 & 0.960 & 0.936 \\
baseline & 1 & 200 & pooled & 0.007 & 0.129 & 0.980 & 0.571 \\
baseline & 1 & 800 & full panel & 0.007 & 0.065 & 0.950 & 0.288 \\
baseline & 1 & 800 & reduced cell & 0.029 & 0.111 & 0.980 & 0.493 \\
baseline & 1 & 800 & pooled & 0.009 & 0.067 & 0.970 & 0.297 \\
baseline & 2 & 200 & full panel & -0.034 & 0.137 & 0.980 & 0.574 \\
baseline & 2 & 200 & reduced cell & -0.020 & 0.187 & 0.980 & 1.032 \\
baseline & 2 & 200 & pooled & -0.035 & 0.140 & 0.970 & 0.590 \\
baseline & 2 & 800 & full panel & 0.013 & 0.065 & 0.970 & 0.295 \\
baseline & 2 & 800 & reduced cell & 0.013 & 0.097 & 1.000 & 0.521 \\
baseline & 2 & 800 & pooled & 0.014 & 0.066 & 0.990 & 0.305 \\
serial & 1 & 200 & full panel & 0.008 & 0.120 & 0.950 & 0.464 \\
serial & 1 & 200 & reduced cell & -0.002 & 0.123 & 0.990 & 0.689 \\
serial & 1 & 200 & pooled & 0.009 & 0.124 & 0.930 & 0.476 \\
serial & 1 & 800 & full panel & -0.005 & 0.058 & 0.950 & 0.234 \\
serial & 1 & 800 & reduced cell & -0.001 & 0.060 & 1.000 & 0.460 \\
serial & 1 & 800 & pooled & -0.004 & 0.059 & 0.970 & 0.242 \\
serial & 2 & 200 & full panel & -0.015 & 0.134 & 0.940 & 0.484 \\
serial & 2 & 200 & reduced cell & -0.023 & 0.126 & 0.990 & 0.734 \\
serial & 2 & 200 & pooled & -0.017 & 0.135 & 0.930 & 0.497 \\
serial & 2 & 800 & full panel & -0.002 & 0.075 & 0.890 & 0.249 \\
serial & 2 & 800 & reduced cell & -0.006 & 0.074 & 1.000 & 0.477 \\
serial & 2 & 800 & pooled & -0.002 & 0.076 & 0.910 & 0.259 \\
staggered & 1 & 200 & full panel & -0.035 & 0.159 & 0.970 & 0.693 \\
staggered & 1 & 200 & reduced cell & 0.012 & 0.240 & 1.000 & 1.265 \\
staggered & 1 & 200 & pooled & 1.417 & 1.424 & 0.000 & 0.514 \\
staggered & 1 & 800 & full panel & -0.007 & 0.097 & 0.940 & 0.389 \\
staggered & 1 & 800 & reduced cell & 0.022 & 0.118 & 1.000 & 0.668 \\
staggered & 1 & 800 & pooled & 1.467 & 1.469 & 0.000 & 0.260 \\
staggered & 3 & 200 & full panel & -0.052 & 0.195 & 0.970 & 0.721 \\
staggered & 3 & 200 & reduced cell & 0.025 & 0.254 & 1.000 & 1.458 \\
staggered & 3 & 200 & pooled & 1.362 & 1.370 & 0.000 & 0.537 \\
staggered & 3 & 800 & full panel & -0.012 & 0.093 & 0.980 & 0.414 \\
staggered & 3 & 800 & reduced cell & 0.019 & 0.134 & 0.990 & 0.732 \\
staggered & 3 & 800 & pooled & 1.425 & 1.427 & 0.000 & 0.275 \\
\end{longtable}
}

{\footnotesize
\begin{longtable}{llrrrrrr}
\caption{ATT in the information-ablation family (100 replications per cell). ``reduced'' uses only the two groups and periods of the cell; ``pooled'' fits one effect surface with \texttt{effect\_by\_cohort=False}.}\label{tab:infoatt}\\
\toprule
Design & $d$ & $N$ & Variant & Bias & RMSE & Cov.95 & Len.95 \\
\midrule
\endfirsthead
\toprule
Design & $d$ & $N$ & Variant & Bias & RMSE & Cov.95 & Len.95 \\
\midrule
\endhead
\bottomrule
\endlastfoot
baseline & 1 & 200 & full panel & 0.004 & 0.109 & 0.960 & 0.426 \\
baseline & 1 & 200 & pooled & 0.006 & 0.106 & 0.970 & 0.427 \\
baseline & 1 & 800 & full panel & 0.007 & 0.054 & 0.940 & 0.215 \\
baseline & 1 & 800 & pooled & 0.007 & 0.054 & 0.940 & 0.215 \\
baseline & 2 & 200 & full panel & -0.029 & 0.129 & 0.940 & 0.438 \\
baseline & 2 & 200 & pooled & -0.030 & 0.129 & 0.940 & 0.440 \\
baseline & 2 & 800 & full panel & 0.014 & 0.058 & 0.920 & 0.226 \\
baseline & 2 & 800 & pooled & 0.015 & 0.057 & 0.930 & 0.227 \\
serial & 1 & 200 & full panel & 0.016 & 0.129 & 0.850 & 0.360 \\
serial & 1 & 200 & pooled & 0.017 & 0.132 & 0.850 & 0.361 \\
serial & 1 & 800 & full panel & -0.005 & 0.062 & 0.880 & 0.183 \\
serial & 1 & 800 & pooled & -0.004 & 0.062 & 0.830 & 0.182 \\
serial & 2 & 200 & full panel & -0.009 & 0.146 & 0.770 & 0.382 \\
serial & 2 & 200 & pooled & -0.008 & 0.149 & 0.770 & 0.383 \\
serial & 2 & 800 & full panel & 0.008 & 0.072 & 0.820 & 0.196 \\
serial & 2 & 800 & pooled & 0.009 & 0.071 & 0.850 & 0.196 \\
staggered & 1 & 200 & full panel & -0.025 & 0.126 & 0.930 & 0.441 \\
staggered & 1 & 200 & pooled & -0.224 & 0.266 & 0.550 & 0.476 \\
staggered & 1 & 800 & full panel & 0.011 & 0.055 & 0.950 & 0.219 \\
staggered & 1 & 800 & pooled & -0.203 & 0.213 & 0.130 & 0.240 \\
staggered & 3 & 200 & full panel & -0.049 & 0.141 & 0.860 & 0.468 \\
staggered & 3 & 200 & pooled & -0.248 & 0.290 & 0.580 & 0.503 \\
staggered & 3 & 800 & full panel & -0.007 & 0.057 & 0.960 & 0.235 \\
staggered & 3 & 800 & pooled & -0.217 & 0.225 & 0.050 & 0.255 \\
\end{longtable}
}

{\footnotesize\setlength{\tabcolsep}{3pt}
\begin{longtable}{lrrrrrrrr}
\caption{Sharp null ($\tau\equiv 0$): rejection rate of the two-sided 5\% posterior test. Published JBES run: 100 replications; new run: available replications.}\label{tab:null}\\
\toprule
Estimator & \multicolumn{4}{c}{$\mathrm{GATT}(4,4)$} & \multicolumn{4}{c}{ATT} \\
\cmidrule(lr){2-5}\cmidrule(lr){6-9}
 & $d1,N200$ & $d2,N200$ & $d1,N800$ & $d2,N800$ & $d1,N200$ & $d2,N200$ & $d1,N800$ & $d2,N800$ \\
\midrule
\endfirsthead
\toprule
Estimator & \multicolumn{4}{c}{$\mathrm{GATT}(4,4)$} & \multicolumn{4}{c}{ATT} \\
\cmidrule(lr){2-5}\cmidrule(lr){6-9}
 & $d1,N200$ & $d2,N200$ & $d1,N800$ & $d2,N800$ & $d1,N200$ & $d2,N200$ & $d1,N800$ & $d2,N800$ \\
\midrule
\endhead
\bottomrule
\endlastfoot
JBES plain (structured) & 0.010 & 0.010 & -- & -- & 0.040 & 0.050 & -- & -- \\
JBES corrected (same-sample) & 0.090 & 0.080 & -- & -- & 0.060 & 0.040 & -- & -- \\
New raw structured & 0.040 & 0.060 & 0.020 & 0.010 & 0.030 & 0.070 & 0.080 & 0.030 \\
New current hybrid, logit & 0.060 & 0.050 & 0.060 & 0.060 & 0.060 & 0.030 & 0.060 & 0.000 \\
New current hybrid, RF & 0.070 & 0.040 & 0.040 & 0.030 & 0.070 & 0.030 & 0.060 & 0.020 \\
New reference fold, logit & 0.060 & 0.060 & 0.090 & 0.070 & 0.100 & 0.100 & 0.070 & 0.050 \\
New reference fold, RF & 0.050 & 0.060 & 0.090 & 0.090 & 0.090 & 0.080 & 0.090 & 0.040 \\
\end{longtable}
}

{\footnotesize
\begin{longtable}{lrrrrrrr}
\caption{Pooled ATT, $N=200$: the newly rerun canonical cells (strong confounder, selection designs, and degree $d=3$) with the published same-sample correction alongside for contrast.}\label{tab:completion_b1}\\
\toprule
Estimator & Bias & SD(err) & RMSE & Cov.95 & Len.95 & $\overline{\mathrm{sd}}/\mathrm{SD}$ & $n$ \\
\midrule
\endfirsthead
Estimator & Bias & SD(err) & RMSE & Cov.95 & Len.95 & $\overline{\mathrm{sd}}/\mathrm{SD}$ & $n$ \\
\midrule
\endhead
\bottomrule
\endlastfoot
\multicolumn{8}{l}{\textit{strong\_confounder}} \\
\multicolumn{8}{l}{\textit{degree $d=1$}} \\
\quad New raw structured & -0.010 & 0.117 & 0.116 & 0.940 & 0.428 & 0.940 & 100 \\
\quad New reference fold, RF & 0.005 & 0.184 & 0.184 & 0.980 & 0.796 & 1.100 & 100 \\
\quad Published corrected (same-sample) & 0.008 & 0.177 & 0.176 & 0.990 & 0.819 & 1.182 & 100 \\
\addlinespace
\multicolumn{8}{l}{\textit{degree $d=2$}} \\
\quad New raw structured & 0.010 & 0.124 & 0.124 & 0.920 & 0.440 & 0.906 & 100 \\
\quad New reference fold, RF & -0.009 & 0.222 & 0.221 & 0.950 & 0.828 & 0.947 & 100 \\
\quad Published corrected (same-sample) & 0.037 & 0.359 & 0.359 & 0.960 & 1.112 & 0.789 & 100 \\
\addlinespace
\multicolumn{8}{l}{\textit{degree $d=3$}} \\
\quad New raw structured & -0.008 & 0.129 & 0.128 & 0.900 & 0.448 & 0.893 & 100 \\
\quad New reference fold, RF & 0.033 & 0.190 & 0.192 & 0.970 & 0.811 & 1.085 & 100 \\
\quad Published corrected (same-sample) & 0.033 & 0.361 & 0.361 & 0.970 & 1.114 & 0.786 & 100 \\
\addlinespace
\multicolumn{8}{l}{\textit{selection\_obs}} \\
\multicolumn{8}{l}{\textit{degree $d=1$}} \\
\quad New raw structured & 0.049 & 0.120 & 0.129 & 0.920 & 0.472 & 1.009 & 100 \\
\quad New reference fold, RF & 0.096 & 0.212 & 0.232 & 0.920 & 0.799 & 0.962 & 100 \\
\quad Published corrected (same-sample) & 0.002 & 0.245 & 0.244 & 0.970 & 1.012 & 1.051 & 100 \\
\addlinespace
\multicolumn{8}{l}{\textit{degree $d=2$}} \\
\quad New raw structured & 0.070 & 0.137 & 0.154 & 0.880 & 0.485 & 0.907 & 100 \\
\quad New reference fold, RF & 0.097 & 0.221 & 0.240 & 0.880 & 0.822 & 0.947 & 100 \\
\quad Published corrected (same-sample) & 0.165 & 1.334 & 1.338 & 0.970 & 2.407 & 0.464 & 100 \\
\addlinespace
\multicolumn{8}{l}{\textit{degree $d=3$}} \\
\quad New raw structured & 0.047 & 0.135 & 0.142 & 0.920 & 0.500 & 0.950 & 100 \\
\quad New reference fold, RF & 0.102 & 0.239 & 0.259 & 0.880 & 0.832 & 0.884 & 100 \\
\quad Published corrected (same-sample) & 0.163 & 1.359 & 1.361 & 0.980 & 2.371 & 0.460 & 100 \\
\addlinespace
\multicolumn{8}{l}{\textit{selection\_both}} \\
\multicolumn{8}{l}{\textit{degree $d=1$}} \\
\quad New raw structured & 0.047 & 0.106 & 0.116 & 0.950 & 0.455 & 1.107 & 100 \\
\quad New reference fold, RF & 0.027 & 0.195 & 0.195 & 0.960 & 0.796 & 1.045 & 100 \\
\quad Published corrected (same-sample) & -0.014 & 0.177 & 0.177 & 0.980 & 0.909 & 1.300 & 100 \\
\addlinespace
\multicolumn{8}{l}{\textit{degree $d=2$}} \\
\quad New raw structured & 0.050 & 0.134 & 0.142 & 0.890 & 0.471 & 0.906 & 100 \\
\quad New reference fold, RF & 0.053 & 0.239 & 0.244 & 0.920 & 0.832 & 0.885 & 100 \\
\quad Published corrected (same-sample) & -0.015 & 0.558 & 0.555 & 1.000 & 1.794 & 0.823 & 100 \\
\addlinespace
\multicolumn{8}{l}{\textit{degree $d=3$}} \\
\quad New raw structured & 0.026 & 0.133 & 0.135 & 0.940 & 0.475 & 0.918 & 100 \\
\quad New reference fold, RF & 0.055 & 0.253 & 0.258 & 0.940 & 0.841 & 0.844 & 100 \\
\quad Published corrected (same-sample) & -0.020 & 0.528 & 0.526 & 0.980 & 1.779 & 0.861 & 100 \\
\addlinespace
\multicolumn{8}{l}{\textit{baseline}} \\
\multicolumn{8}{l}{\textit{degree $d=3$}} \\
\quad New raw structured & 0.006 & 0.116 & 0.115 & 0.940 & 0.445 & 0.985 & 100 \\
\quad New reference fold, RF & 0.009 & 0.220 & 0.219 & 0.940 & 0.806 & 0.935 & 100 \\
\quad Published corrected (same-sample) & 0.112 & 1.096 & 1.096 & 0.970 & 1.520 & 0.353 & 100 \\
\addlinespace
\multicolumn{8}{l}{\textit{serial}} \\
\multicolumn{8}{l}{\textit{degree $d=3$}} \\
\quad New raw structured & 0.015 & 0.172 & 0.172 & 0.710 & 0.382 & 0.568 & 100 \\
\quad New reference fold, RF & 0.027 & 0.197 & 0.198 & 0.900 & 0.696 & 0.901 & 100 \\
\quad Published corrected (same-sample) & 0.078 & 0.599 & 0.601 & 0.960 & 1.268 & 0.538 & 100 \\
\addlinespace
\multicolumn{8}{l}{\textit{null}} \\
\multicolumn{8}{l}{\textit{degree $d=3$}} \\
\quad New raw structured & -0.001 & 0.122 & 0.122 & 0.945 & 0.438 & 0.919 & 200 \\
\quad New reference fold, RF & -0.012 & 0.232 & 0.231 & 0.870 & 0.724 & 0.795 & 200 \\
\quad Published corrected (same-sample) & 0.123 & 1.092 & 1.094 & 0.960 & 1.409 & 0.331 & 100 \\
\addlinespace
\end{longtable}
}

{\footnotesize
\begin{longtable}{lrrrrrrr}
\caption{$\mathrm{GATT}(g{=}4,t{=}4)$ sweeps: every cell the correction-audit family skipped, with the published same-sample correction alongside for contrast.}\label{tab:completion_sweep}\\
\toprule
Estimator & Bias & SD(err) & RMSE & Cov.95 & Len.95 & $\overline{\mathrm{sd}}/\mathrm{SD}$ & $n$ \\
\midrule
\endfirsthead
Estimator & Bias & SD(err) & RMSE & Cov.95 & Len.95 & $\overline{\mathrm{sd}}/\mathrm{SD}$ & $n$ \\
\midrule
\endhead
\bottomrule
\endlastfoot
\multicolumn{8}{l}{\textit{baseline\_sweep}} \\
\multicolumn{8}{l}{\textit{$N=50$, degree $d=1$}} \\
\quad New raw structured & -0.013 & 0.229 & 0.228 & 0.980 & 1.057 & 1.180 & 100 \\
\quad New reference fold, RF & -0.124 & 0.520 & 0.532 & 0.960 & 2.144 & 1.041 & 100 \\
\quad Published corrected (same-sample) & 0.083 & 0.705 & 0.707 & 0.940 & 2.889 & 1.041 & 100 \\
\addlinespace
\multicolumn{8}{l}{\textit{$N=100$, degree $d=1$}} \\
\quad New raw structured & -0.019 & 0.215 & 0.214 & 0.940 & 0.776 & 0.922 & 100 \\
\quad New reference fold, RF & 0.023 & 0.412 & 0.410 & 0.920 & 1.461 & 0.898 & 100 \\
\quad Published corrected (same-sample) & 0.019 & 0.343 & 0.342 & 0.950 & 1.567 & 1.157 & 100 \\
\addlinespace
\multicolumn{8}{l}{\textit{$N=400$, degree $d=1$}} \\
\quad New raw structured & 0.012 & 0.091 & 0.091 & 0.980 & 0.401 & 1.132 & 100 \\
\quad New reference fold, RF & 0.000 & 0.190 & 0.189 & 0.920 & 0.683 & 0.919 & 100 \\
\quad Published corrected (same-sample) & 0.001 & 0.154 & 0.154 & 0.950 & 0.652 & 1.082 & 100 \\
\addlinespace
\multicolumn{8}{l}{\textit{$N=50$, degree $d=2$}} \\
\quad New raw structured & 0.019 & 0.267 & 0.266 & 0.950 & 1.085 & 1.040 & 100 \\
\quad New reference fold, RF & 0.061 & 0.685 & 0.685 & 0.920 & 2.245 & 0.827 & 100 \\
\quad Published corrected (same-sample) & 0.186 & 2.523 & 2.517 & 0.970 & 4.778 & 0.483 & 100 \\
\addlinespace
\multicolumn{8}{l}{\textit{$N=100$, degree $d=2$}} \\
\quad New raw structured & 0.005 & 0.187 & 0.186 & 0.960 & 0.795 & 1.088 & 100 \\
\quad New reference fold, RF & 0.001 & 0.381 & 0.379 & 0.960 & 1.446 & 0.963 & 100 \\
\quad Published corrected (same-sample) & 0.081 & 1.259 & 1.255 & 0.940 & 2.963 & 0.598 & 100 \\
\addlinespace
\multicolumn{8}{l}{\textit{$N=400$, degree $d=2$}} \\
\quad New raw structured & 0.012 & 0.102 & 0.102 & 0.960 & 0.413 & 1.034 & 100 \\
\quad New reference fold, RF & -0.002 & 0.148 & 0.147 & 0.980 & 0.690 & 1.192 & 100 \\
\quad Published corrected (same-sample) & -0.200 & 1.310 & 1.319 & 0.960 & 1.482 & 0.298 & 100 \\
\addlinespace
\multicolumn{8}{l}{\textit{baseline\_sweep\_d3}} \\
\multicolumn{8}{l}{\textit{$N=50$, degree $d=3$}} \\
\quad New raw structured & 0.022 & 0.256 & 0.256 & 0.980 & 1.122 & 1.117 & 100 \\
\quad New reference fold, RF & 0.041 & 0.573 & 0.571 & 0.930 & 2.154 & 0.955 & 100 \\
\quad Published corrected (same-sample) & 0.179 & 2.461 & 2.455 & 0.970 & 4.810 & 0.504 & 100 \\
\addlinespace
\multicolumn{8}{l}{\textit{$N=100$, degree $d=3$}} \\
\quad New raw structured & 0.003 & 0.219 & 0.217 & 0.930 & 0.802 & 0.938 & 100 \\
\quad New reference fold, RF & -0.025 & 0.336 & 0.335 & 0.980 & 1.441 & 1.091 & 100 \\
\quad Published corrected (same-sample) & 0.058 & 1.151 & 1.147 & 0.960 & 2.985 & 0.675 & 100 \\
\addlinespace
\multicolumn{8}{l}{\textit{$N=400$, degree $d=3$}} \\
\quad New raw structured & 0.003 & 0.102 & 0.102 & 0.960 & 0.419 & 1.049 & 100 \\
\quad New reference fold, RF & 0.019 & 0.169 & 0.169 & 0.980 & 0.690 & 1.041 & 100 \\
\quad Published corrected (same-sample) & -0.194 & 1.249 & 1.257 & 0.960 & 1.455 & 0.304 & 100 \\
\addlinespace
\multicolumn{8}{l}{\textit{$N=800$, degree $d=3$}} \\
\quad New raw structured & 0.003 & 0.074 & 0.074 & 0.970 & 0.304 & 1.053 & 100 \\
\quad New reference fold, RF & -0.001 & 0.136 & 0.135 & 0.900 & 0.484 & 0.907 & 100 \\
\quad Published corrected (same-sample) & -0.049 & 0.337 & 0.339 & 0.980 & 0.643 & 0.492 & 100 \\
\addlinespace
\multicolumn{8}{l}{\textit{serial\_sweep}} \\
\multicolumn{8}{l}{\textit{$N=50$, degree $d=1$}} \\
\quad New raw structured & 0.054 & 0.308 & 0.311 & 0.840 & 0.909 & 0.755 & 100 \\
\quad New reference fold, RF & 0.020 & 0.337 & 0.335 & 0.950 & 1.422 & 1.074 & 100 \\
\quad Published corrected (same-sample) & 0.030 & 0.479 & 0.478 & 0.970 & 1.955 & 1.039 & 100 \\
\addlinespace
\multicolumn{8}{l}{\textit{$N=100$, degree $d=1$}} \\
\quad New raw structured & -0.003 & 0.200 & 0.199 & 0.900 & 0.645 & 0.825 & 100 \\
\quad New reference fold, RF & 0.003 & 0.235 & 0.234 & 0.950 & 0.979 & 1.061 & 100 \\
\quad Published corrected (same-sample) & 0.014 & 0.201 & 0.201 & 0.990 & 1.098 & 1.383 & 100 \\
\addlinespace
\multicolumn{8}{l}{\textit{$N=400$, degree $d=1$}} \\
\quad New raw structured & 0.009 & 0.099 & 0.099 & 0.910 & 0.329 & 0.847 & 100 \\
\quad New reference fold, RF & -0.002 & 0.125 & 0.125 & 0.990 & 0.491 & 1.000 & 100 \\
\quad Published corrected (same-sample) & 0.001 & 0.091 & 0.090 & 1.000 & 0.470 & 1.331 & 100 \\
\addlinespace
\multicolumn{8}{l}{\textit{$N=50$, degree $d=2$}} \\
\quad New raw structured & 0.026 & 0.303 & 0.302 & 0.890 & 0.950 & 0.802 & 100 \\
\quad New reference fold, RF & 0.048 & 0.398 & 0.399 & 0.940 & 1.441 & 0.917 & 100 \\
\quad Published corrected (same-sample) & 0.211 & 1.860 & 1.863 & 1.000 & 3.596 & 0.506 & 100 \\
\addlinespace
\multicolumn{8}{l}{\textit{$N=100$, degree $d=2$}} \\
\quad New raw structured & 0.005 & 0.187 & 0.186 & 0.920 & 0.670 & 0.912 & 100 \\
\quad New reference fold, RF & -0.005 & 0.222 & 0.220 & 0.950 & 0.989 & 1.134 & 100 \\
\quad Published corrected (same-sample) & 0.085 & 0.886 & 0.886 & 0.990 & 2.283 & 0.657 & 100 \\
\addlinespace
\multicolumn{8}{l}{\textit{$N=400$, degree $d=2$}} \\
\quad New raw structured & -0.012 & 0.101 & 0.101 & 0.900 & 0.346 & 0.879 & 100 \\
\quad New reference fold, RF & 0.001 & 0.106 & 0.106 & 0.970 & 0.496 & 1.192 & 100 \\
\quad Published corrected (same-sample) & -0.113 & 0.835 & 0.839 & 0.990 & 1.018 & 0.316 & 100 \\
\addlinespace
\multicolumn{8}{l}{\textit{serial\_sweep\_d3}} \\
\multicolumn{8}{l}{\textit{$N=50$, degree $d=3$}} \\
\quad New raw structured & -0.030 & 0.286 & 0.287 & 0.930 & 0.951 & 0.847 & 100 \\
\quad New reference fold, RF & -0.036 & 0.344 & 0.344 & 0.940 & 1.428 & 1.053 & 100 \\
\quad Published corrected (same-sample) & 0.240 & 1.604 & 1.614 & 0.980 & 3.592 & 0.581 & 100 \\
\addlinespace
\multicolumn{8}{l}{\textit{$N=100$, degree $d=3$}} \\
\quad New raw structured & -0.020 & 0.174 & 0.174 & 0.980 & 0.687 & 1.009 & 100 \\
\quad New reference fold, RF & 0.018 & 0.233 & 0.232 & 0.960 & 0.995 & 1.088 & 100 \\
\quad Published corrected (same-sample) & 0.065 & 0.551 & 0.552 & 0.980 & 2.155 & 1.002 & 100 \\
\addlinespace
\multicolumn{8}{l}{\textit{$N=400$, degree $d=3$}} \\
\quad New raw structured & -0.005 & 0.097 & 0.096 & 0.940 & 0.352 & 0.926 & 100 \\
\quad New reference fold, RF & -0.007 & 0.114 & 0.113 & 0.950 & 0.490 & 1.102 & 100 \\
\quad Published corrected (same-sample) & -0.094 & 0.656 & 0.660 & 1.000 & 0.922 & 0.359 & 100 \\
\addlinespace
\multicolumn{8}{l}{\textit{$N=800$, degree $d=3$}} \\
\quad New raw structured & -0.007 & 0.058 & 0.058 & 0.960 & 0.251 & 1.125 & 100 \\
\quad New reference fold, RF & 0.005 & 0.073 & 0.072 & 0.990 & 0.344 & 1.210 & 100 \\
\quad Published corrected (same-sample) & -0.018 & 0.145 & 0.146 & 0.990 & 0.442 & 0.773 & 100 \\
\addlinespace
\end{longtable}
}

{\footnotesize
\begin{longtable}{lrrrrrrr}
\caption{Staggered event-study and pooled ATT cells at all three linearity degrees ($N=200$), with the published same-sample correction alongside for contrast.}\label{tab:completion_staggered}\\
\toprule
Estimator & Bias & SD(err) & RMSE & Cov.95 & Len.95 & $\overline{\mathrm{sd}}/\mathrm{SD}$ & $n$ \\
\midrule
\endfirsthead
Estimator & Bias & SD(err) & RMSE & Cov.95 & Len.95 & $\overline{\mathrm{sd}}/\mathrm{SD}$ & $n$ \\
\midrule
\endhead
\bottomrule
\endlastfoot
\multicolumn{8}{l}{\textit{degree $d=1$}} \\
\multicolumn{8}{l}{\textit{ES, k$=$0}} \\
\quad New raw structured & 0.107 & 0.109 & 0.153 & 0.840 & 0.460 & 1.081 & 100 \\
\quad New reference fold, RF & -0.046 & 0.168 & 0.174 & 0.980 & 0.808 & 1.224 & 100 \\
\quad Published corrected (same-sample) & -0.011 & 0.160 & 0.159 & 1.000 & 0.849 & 1.356 & 100 \\
\addlinespace
\multicolumn{8}{l}{\textit{ES, k$=$1}} \\
\quad New raw structured & -0.006 & 0.143 & 0.142 & 0.930 & 0.528 & 0.949 & 100 \\
\quad New reference fold, RF & -0.023 & 0.211 & 0.211 & 0.970 & 0.955 & 1.154 & 100 \\
\quad Published corrected (same-sample) & -0.015 & 0.192 & 0.191 & 0.990 & 1.012 & 1.343 & 100 \\
\addlinespace
\multicolumn{8}{l}{\textit{ES, k$=$2}} \\
\quad New raw structured & -0.070 & 0.176 & 0.189 & 0.940 & 0.669 & 0.973 & 100 \\
\quad New reference fold, RF & -0.001 & 0.302 & 0.301 & 0.960 & 1.218 & 1.026 & 100 \\
\quad Published corrected (same-sample) & -0.003 & 0.247 & 0.245 & 1.000 & 1.279 & 1.321 & 100 \\
\addlinespace
\multicolumn{8}{l}{\textit{ES, k$=$3}} \\
\quad New raw structured & -0.355 & 0.259 & 0.439 & 0.660 & 0.901 & 0.897 & 100 \\
\quad New reference fold, RF & -0.031 & 0.382 & 0.381 & 0.960 & 1.665 & 1.112 & 100 \\
\quad Published corrected (same-sample) & -0.004 & 0.395 & 0.393 & 0.970 & 1.748 & 1.126 & 100 \\
\addlinespace
\multicolumn{8}{l}{\textit{ATT, ATT}} \\
\quad New raw structured & -0.023 & 0.119 & 0.120 & 0.940 & 0.445 & 0.963 & 100 \\
\quad New reference fold, RF & -0.027 & 0.202 & 0.203 & 0.970 & 0.860 & 1.084 & 100 \\
\quad Published corrected (same-sample) & -0.010 & 0.178 & 0.177 & 1.000 & 0.914 & 1.310 & 100 \\
\addlinespace
\multicolumn{8}{l}{\textit{degree $d=2$}} \\
\multicolumn{8}{l}{\textit{ES, k$=$0}} \\
\quad New raw structured & 0.098 & 0.116 & 0.152 & 0.870 & 0.470 & 1.042 & 100 \\
\quad New reference fold, RF & -0.008 & 0.160 & 0.159 & 0.980 & 0.792 & 1.266 & 100 \\
\quad Published corrected (same-sample) & 0.191 & 2.060 & 2.059 & 1.000 & 2.356 & 0.292 & 100 \\
\addlinespace
\multicolumn{8}{l}{\textit{ES, k$=$1}} \\
\quad New raw structured & -0.006 & 0.151 & 0.150 & 0.940 & 0.540 & 0.916 & 100 \\
\quad New reference fold, RF & 0.033 & 0.209 & 0.211 & 0.970 & 0.945 & 1.153 & 100 \\
\quad Published corrected (same-sample) & 0.178 & 1.822 & 1.822 & 0.990 & 2.370 & 0.336 & 100 \\
\addlinespace
\multicolumn{8}{l}{\textit{ES, k$=$2}} \\
\quad New raw structured & -0.049 & 0.177 & 0.183 & 0.930 & 0.684 & 0.987 & 100 \\
\quad New reference fold, RF & 0.045 & 0.255 & 0.258 & 0.990 & 1.211 & 1.211 & 100 \\
\quad Published corrected (same-sample) & 0.345 & 3.627 & 3.626 & 0.990 & 3.809 & 0.273 & 100 \\
\addlinespace
\multicolumn{8}{l}{\textit{ES, k$=$3}} \\
\quad New raw structured & -0.360 & 0.292 & 0.463 & 0.660 & 0.921 & 0.808 & 100 \\
\quad New reference fold, RF & 0.033 & 0.438 & 0.437 & 0.950 & 1.669 & 0.972 & 100 \\
\quad Published corrected (same-sample) & 0.267 & 3.368 & 3.362 & 0.980 & 3.737 & 0.286 & 100 \\
\addlinespace
\multicolumn{8}{l}{\textit{ATT, ATT}} \\
\quad New raw structured & -0.023 & 0.127 & 0.129 & 0.920 & 0.456 & 0.920 & 100 \\
\quad New reference fold, RF & 0.022 & 0.188 & 0.189 & 0.990 & 0.858 & 1.157 & 100 \\
\quad Published corrected (same-sample) & 0.231 & 2.424 & 2.423 & 1.000 & 2.491 & 0.266 & 100 \\
\addlinespace
\multicolumn{8}{l}{\textit{degree $d=3$}} \\
\multicolumn{8}{l}{\textit{ES, k$=$0}} \\
\quad New raw structured & 0.088 & 0.112 & 0.142 & 0.930 & 0.478 & 1.093 & 100 \\
\quad New reference fold, RF & -0.030 & 0.170 & 0.171 & 1.000 & 0.786 & 1.180 & 100 \\
\quad Published corrected (same-sample) & 0.238 & 2.548 & 2.547 & 1.000 & 2.449 & 0.249 & 100 \\
\addlinespace
\multicolumn{8}{l}{\textit{ES, k$=$1}} \\
\quad New raw structured & -0.053 & 0.142 & 0.151 & 0.930 & 0.558 & 1.004 & 100 \\
\quad New reference fold, RF & 0.002 & 0.201 & 0.200 & 0.970 & 0.923 & 1.170 & 100 \\
\quad Published corrected (same-sample) & 0.103 & 1.373 & 1.370 & 1.000 & 2.241 & 0.424 & 100 \\
\addlinespace
\multicolumn{8}{l}{\textit{ES, k$=$2}} \\
\quad New raw structured & -0.140 & 0.181 & 0.229 & 0.890 & 0.700 & 0.988 & 100 \\
\quad New reference fold, RF & -0.019 & 0.272 & 0.271 & 0.970 & 1.186 & 1.109 & 100 \\
\quad Published corrected (same-sample) & 0.247 & 3.204 & 3.197 & 0.990 & 3.608 & 0.294 & 100 \\
\addlinespace
\multicolumn{8}{l}{\textit{ES, k$=$3}} \\
\quad New raw structured & -0.428 & 0.259 & 0.500 & 0.580 & 0.935 & 0.924 & 100 \\
\quad New reference fold, RF & -0.024 & 0.402 & 0.401 & 0.960 & 1.621 & 1.024 & 100 \\
\quad Published corrected (same-sample) & 0.155 & 2.603 & 2.595 & 0.990 & 3.683 & 0.359 & 100 \\
\addlinespace
\multicolumn{8}{l}{\textit{ATT, ATT}} \\
\quad New raw structured & -0.071 & 0.119 & 0.138 & 0.870 & 0.468 & 1.011 & 100 \\
\quad New reference fold, RF & -0.016 & 0.190 & 0.189 & 0.950 & 0.835 & 1.121 & 100 \\
\quad Published corrected (same-sample) & 0.185 & 2.235 & 2.231 & 1.000 & 2.415 & 0.280 & 100 \\
\addlinespace
\end{longtable}
}

{\footnotesize
\begin{longtable}{lrrrrrrr}
\caption{Sharp-null $d=3$ rerun ($N=200$, 200 replications): bias, dispersion, coverage, and the two-sided 5\% rejection rate with its Monte-Carlo standard error.}\label{tab:completion_null}\\
\toprule
Estimator & Bias & SD(err) & RMSE & Cov.95 & Size (5\%) & MCSE & $n$ \\
\midrule
\endfirsthead
Estimator & Bias & SD(err) & RMSE & Cov.95 & Size (5\%) & MCSE & $n$ \\
\midrule
\endhead
\bottomrule
\endlastfoot
\multicolumn{8}{l}{\textit{ATT, ATT}} \\
\quad New raw structured & -0.001 & 0.122 & 0.122 & 0.945 & 0.050 & 0.015 & 200 \\
\quad New reference fold, RF & -0.012 & 0.232 & 0.231 & 0.870 & 0.130 & 0.024 & 200 \\
\addlinespace
\multicolumn{8}{l}{\textit{ES, k$=$0}} \\
\quad New raw structured & -0.001 & 0.132 & 0.132 & 0.970 & 0.030 & 0.012 & 200 \\
\quad New reference fold, RF & -0.009 & 0.264 & 0.264 & 0.900 & 0.100 & 0.021 & 200 \\
\addlinespace
\multicolumn{8}{l}{\textit{GATT, g$=$4\_t$=$4}} \\
\quad New raw structured & -0.001 & 0.132 & 0.132 & 0.970 & 0.030 & 0.012 & 200 \\
\quad New reference fold, RF & -0.009 & 0.264 & 0.264 & 0.900 & 0.100 & 0.021 & 200 \\
\addlinespace
\end{longtable}
}

{\footnotesize
\begin{longtable}{lrrrrrrr}
\caption{Pre-trend diagnostic operating characteristics at degree $d=1$: slope-rule, $X_1$-contrast, and any-$k$ detection rates with slope bias, RMSE, and coverage. The fold-pooled diagnostic is an exploratory aggregation, not a theorem-covered estimator.}\label{tab:completion_pt}\\
\toprule
Diagnostic & Slope & Contrast & Any-$k$ & Bias & RMSE & Cov.95 & $n$ \\
\midrule
\endfirsthead
Diagnostic & Slope & Contrast & Any-$k$ & Bias & RMSE & Cov.95 & $n$ \\
\midrule
\endhead
\bottomrule
\endlastfoot
\multicolumn{8}{l}{\textit{PT\_hold}} \\
\quad New raw diagnostic & 0.015 & 0.000 & 0.005 & 0.031 & 0.061 & 0.985 & 200 \\
\quad New reference-fold diagnostic & 0.105 & 0.000 & 0.040 & 0.061 & 0.078 & 0.890 & 200 \\
\addlinespace
\multicolumn{8}{l}{\textit{PT\_conditional}} \\
\quad New raw diagnostic & 0.060 & 0.000 & 0.015 & 0.061 & 0.083 & 0.940 & 200 \\
\quad New reference-fold diagnostic & 0.275 & 0.000 & 0.120 & 0.102 & 0.115 & 0.725 & 200 \\
\addlinespace
\multicolumn{8}{l}{\textit{PT\_violation\_g05}} \\
\quad New raw diagnostic & 0.130 & 0.000 & 0.055 & 0.025 & 0.061 & 0.965 & 200 \\
\quad New reference-fold diagnostic & 0.375 & 0.005 & 0.180 & 0.052 & 0.074 & 0.925 & 200 \\
\addlinespace
\multicolumn{8}{l}{\textit{PT\_violation\_g10}} \\
\quad New raw diagnostic & 0.385 & 0.000 & 0.240 & 0.014 & 0.067 & 0.955 & 200 \\
\quad New reference-fold diagnostic & 0.600 & 0.000 & 0.410 & 0.037 & 0.070 & 0.930 & 200 \\
\addlinespace
\multicolumn{8}{l}{\textit{PT\_violation\_g20}} \\
\quad New raw diagnostic & 0.855 & 0.000 & 0.755 & 0.007 & 0.069 & 0.945 & 200 \\
\quad New reference-fold diagnostic & 0.940 & 0.010 & 0.930 & 0.026 & 0.069 & 0.935 & 200 \\
\addlinespace
\multicolumn{8}{l}{\textit{PT\_violation\_g40}} \\
\quad New raw diagnostic & 1.000 & 0.000 & 1.000 & -0.003 & 0.070 & 0.955 & 200 \\
\quad New reference-fold diagnostic & 1.000 & 0.010 & 1.000 & 0.003 & 0.067 & 0.960 & 200 \\
\addlinespace
\multicolumn{8}{l}{\textit{PT\_violation\_a20}} \\
\quad New raw diagnostic & 0.870 & 0.000 & 0.725 & 0.009 & 0.068 & 0.970 & 200 \\
\quad New reference-fold diagnostic & 0.945 & 0.005 & 0.890 & 0.024 & 0.067 & 0.965 & 200 \\
\addlinespace
\multicolumn{8}{l}{\textit{PT\_violation\_het10}} \\
\quad New raw diagnostic & 0.015 & 0.020 & 0.000 & 0.030 & 0.061 & 0.975 & 200 \\
\quad New reference-fold diagnostic & 0.100 & 0.025 & 0.035 & 0.060 & 0.078 & 0.895 & 200 \\
\addlinespace
\multicolumn{8}{l}{\textit{PT\_violation\_het20}} \\
\quad New raw diagnostic & 0.060 & 0.450 & 0.035 & 0.038 & 0.074 & 0.930 & 200 \\
\quad New reference-fold diagnostic & 0.190 & 0.365 & 0.095 & 0.067 & 0.088 & 0.840 & 200 \\
\addlinespace
\multicolumn{8}{l}{\textit{PT\_violation\_het40}} \\
\quad New raw diagnostic & 0.095 & 1.000 & 0.040 & 0.036 & 0.076 & 0.940 & 200 \\
\quad New reference-fold diagnostic & 0.185 & 1.000 & 0.110 & 0.068 & 0.092 & 0.840 & 200 \\
\addlinespace
\end{longtable}
}

{\footnotesize
\begin{longtable}{lrrrrrrr}
\caption{Pre-trend diagnostic operating characteristics at degree $d=3$: slope-rule, $X_1$-contrast, and any-$k$ detection rates with slope bias, RMSE, and coverage. The fold-pooled diagnostic is an exploratory aggregation, not a theorem-covered estimator.}\label{tab:completion_pt_d3}\\
\toprule
Diagnostic & Slope & Contrast & Any-$k$ & Bias & RMSE & Cov.95 & $n$ \\
\midrule
\endfirsthead
Diagnostic & Slope & Contrast & Any-$k$ & Bias & RMSE & Cov.95 & $n$ \\
\midrule
\endhead
\bottomrule
\endlastfoot
\multicolumn{8}{l}{\textit{PT\_hold}} \\
\quad New raw diagnostic & 0.040 & 0.000 & 0.015 & 0.037 & 0.071 & 0.960 & 200 \\
\quad New reference-fold diagnostic & 0.135 & 0.015 & 0.050 & 0.067 & 0.088 & 0.865 & 200 \\
\addlinespace
\multicolumn{8}{l}{\textit{PT\_conditional}} \\
\quad New raw diagnostic & 0.070 & 0.000 & 0.045 & 0.064 & 0.092 & 0.930 & 200 \\
\quad New reference-fold diagnostic & 0.295 & 0.010 & 0.155 & 0.112 & 0.128 & 0.705 & 200 \\
\addlinespace
\multicolumn{8}{l}{\textit{PT\_violation\_g05}} \\
\quad New raw diagnostic & 0.160 & 0.000 & 0.070 & 0.029 & 0.067 & 0.955 & 200 \\
\quad New reference-fold diagnostic & 0.360 & 0.000 & 0.200 & 0.056 & 0.080 & 0.910 & 200 \\
\addlinespace
\multicolumn{8}{l}{\textit{PT\_violation\_g10}} \\
\quad New raw diagnostic & 0.385 & 0.000 & 0.215 & 0.024 & 0.069 & 0.955 & 200 \\
\quad New reference-fold diagnostic & 0.650 & 0.005 & 0.480 & 0.048 & 0.076 & 0.920 & 200 \\
\addlinespace
\multicolumn{8}{l}{\textit{PT\_violation\_g20}} \\
\quad New raw diagnostic & 0.890 & 0.000 & 0.760 & 0.018 & 0.069 & 0.945 & 200 \\
\quad New reference-fold diagnostic & 0.975 & 0.000 & 0.935 & 0.036 & 0.070 & 0.940 & 200 \\
\addlinespace
\multicolumn{8}{l}{\textit{PT\_violation\_g40}} \\
\quad New raw diagnostic & 1.000 & 0.000 & 1.000 & -0.004 & 0.060 & 0.980 & 200 \\
\quad New reference-fold diagnostic & 1.000 & 0.010 & 1.000 & 0.003 & 0.057 & 0.990 & 200 \\
\addlinespace
\multicolumn{8}{l}{\textit{PT\_violation\_a20}} \\
\quad New raw diagnostic & 0.865 & 0.000 & 0.740 & 0.016 & 0.063 & 0.980 & 200 \\
\quad New reference-fold diagnostic & 0.970 & 0.015 & 0.920 & 0.034 & 0.066 & 0.975 & 200 \\
\addlinespace
\multicolumn{8}{l}{\textit{PT\_violation\_het10}} \\
\quad New raw diagnostic & 0.070 & 0.035 & 0.025 & 0.040 & 0.073 & 0.920 & 200 \\
\quad New reference-fold diagnostic & 0.150 & 0.035 & 0.090 & 0.071 & 0.090 & 0.865 & 200 \\
\addlinespace
\multicolumn{8}{l}{\textit{PT\_violation\_het20}} \\
\quad New raw diagnostic & 0.070 & 0.345 & 0.030 & 0.042 & 0.072 & 0.935 & 200 \\
\quad New reference-fold diagnostic & 0.125 & 0.275 & 0.075 & 0.073 & 0.089 & 0.820 & 200 \\
\addlinespace
\multicolumn{8}{l}{\textit{PT\_violation\_het40}} \\
\quad New raw diagnostic & 0.135 & 1.000 & 0.095 & 0.046 & 0.084 & 0.900 & 200 \\
\quad New reference-fold diagnostic & 0.260 & 1.000 & 0.160 & 0.078 & 0.102 & 0.810 & 200 \\
\addlinespace
\end{longtable}
}


\begin{table}[htbp]
\centering
\footnotesize
\caption{Head-to-head summary of the old and new corrections on
$\mathrm{GATT}(4,4)$: RMSE with 95\% coverage in parentheses. The theory-aligned
reference fold uses the random-forest pilot; ``clip .05'' is the same-sample logit with
the clipping constant set to $0.05$.}
\label{tab:headtohead}
\begin{tabular}{lrrrr}
\toprule
Cell & Old corrected & Same-sample logit & Same-sample clip .05 & Reference fold RF \\
\midrule
baseline, $d=1$, $N=200$ & 0.231 (0.95) & 0.214 (0.96) & 0.214 (0.96) & 0.243 (0.91) \\
baseline, $d=2$, $N=200$ & 1.040 (0.96) & 5.947 (0.98) & 0.258 (0.98) & 0.234 (0.97) \\
baseline, $d=2$, $N=800$ & 0.318 (0.99) & 0.102 (0.96) & 0.102 (0.96) & 0.114 (0.95) \\
serial, $d=1$, $N=200$ & 0.136 (0.99) & 0.145 (0.98) & 0.145 (0.98) & 0.160 (0.96) \\
serial, $d=2$, $N=200$ & 0.269 (0.99) & 0.264 (0.99) & 0.158 (0.99) & 0.164 (0.98) \\
\bottomrule
\end{tabular}
\end{table}

\begin{table}[htbp]
\centering
\footnotesize
\caption{Head-to-head summary against the published benchmarks on
$\mathrm{GATT}(4,4)$: RMSE with 95\% coverage in parentheses. DiD-BCF entries are the
new runs; benchmark entries are the published JBES runs.}
\label{tab:benchhead}
\begin{tabular}{lrrrrr}
\toprule
Cell & Raw & Reference fold RF & Callaway--Sant'Anna & DoubleML & grf-DiD \\
\midrule
baseline, $d=1$, $N=200$ & 0.131 (0.98) & 0.243 (0.91) & 0.221 (0.97) & 0.304 (0.92) & 0.222 (0.81) \\
baseline, $d=2$, $N=200$ & 0.135 (0.96) & 0.234 (0.97) & 0.225 (0.96) & 0.280 (0.94) & 0.223 (0.81) \\
baseline, $d=1$, $N=800$ & 0.069 (0.96) & 0.130 (0.92) & 0.107 (0.98) & 0.114 (0.97) & 0.111 (0.81) \\
serial, $d=1$, $N=200$ & 0.126 (0.92) & 0.160 (0.96) & 0.129 (0.99) & 0.165 (0.99) & 0.133 (0.91) \\
serial, $d=2$, $N=200$ & 0.141 (0.88) & 0.164 (0.98) & 0.135 (0.99) & 0.170 (0.97) & 0.133 (0.90) \\
\bottomrule
\end{tabular}
\end{table}

\begin{thebibliography}{9}
\bibitem{souto2025}
H.~G. Souto and F.~Louzada.
\newblock Forests for Differences: Robust Causal Inference Beyond Parametric DiD.
\newblock arXiv:2505.09706, 2025.
\newblock \url{https://doi.org/10.48550/arXiv.2505.09706}.

\bibitem{santanna2020}
P.~H.~C. Sant'Anna and J.~Zhao.
\newblock Doubly robust difference-in-differences estimators.
\newblock \emph{Journal of Econometrics}, 219(1):101--122, 2020.
\newblock \url{https://doi.org/10.1016/j.jeconom.2020.05.003}.

\bibitem{callaway2021}
B.~Callaway and P.~H.~C. Sant'Anna.
\newblock Difference-in-Differences with multiple time periods.
\newblock \emph{Journal of Econometrics}, 225(2):200--230, 2021.
\newblock \url{https://doi.org/10.1016/j.jeconom.2020.12.001}.
\end{thebibliography}

\end{document}
